
\Data\ID{1003034101}
\Updated{2020-07-15 03:07:12}
\Level{A}
\SubArea{Elementary arithmetics}
\LANGen{
\Question{
Evaluating the expression \( \frac12 + \frac23 \) we get:}
{\( \frac76 \)}{\( \frac35 \)}{\( \frac7{12} \)}{\( \frac36 \)}{}{}}
\LANGpl{
\Question{
Wartością wyrażenia \( \frac12 + \frac23 \) jest:}
{\( \frac76 \)}{\( \frac35 \)}{\( \frac7{12} \)}{\( \frac36 \)}{}{}}
\LANGcs{
\Question{
Hodnota výrazu  \( \frac12 + \frac23 \) je:}
{\( \frac76 \)}{\( \frac35 \)}{\( \frac7{12} \)}{\( \frac36 \)}{}{}}
\LANGsk{
\Question{
Hodnota výrazu \( \frac12 + \frac23 \) je:}
{\( \frac76 \)}{\( \frac35 \)}{\( \frac7{12} \)}{\( \frac36 \)}{}{}}
\LANGes{
\Question{
Al calcular la expresión \( \frac12 + \frac23 \) obtenemos:}
{\( \frac76 \)}{\( \frac35 \)}{\( \frac7{12} \)}{\( \frac36 \)}{}{}}
\End

\Data\ID{1003034102}
\Updated{2020-07-15 03:07:12}
\Level{A}
\SubArea{Elementary arithmetics}
\LANGen{
\Question{
Evaluating the expression \( \frac27 \cdot 3 \) we get:}
{\( \frac67 \)}{\( \frac57 \)}{\( \frac6{21} \)}{\( \frac2{21} \)}{}{}}
\LANGpl{
\Question{
Wartością wyrażenia \( \frac27 \cdot 3 \) jest:}
{\( \frac67 \)}{\( \frac57 \)}{\( \frac6{21} \)}{\( \frac2{21} \)}{}{}}
\LANGcs{
\Question{
Hodnota výrazu \( \frac27 \cdot 3 \) je:}
{\( \frac67 \)}{\( \frac57 \)}{\( \frac6{21} \)}{\( \frac2{21} \)}{}{}}
\LANGsk{
\Question{
Hodnota výrazu \( \frac27 \cdot 3 \) je:}
{\( \frac67 \)}{\( \frac57 \)}{\( \frac6{21} \)}{\( \frac2{21} \)}{}{}}
\LANGes{
\Question{
Al evaluar la expresión  \( \frac27 \cdot 3 \) obtenemos:}
{\( \frac67 \)}{\( \frac57 \)}{\( \frac6{21} \)}{\( \frac2{21} \)}{}{}}
\End

\Data\ID{1003034103}
\Updated{2020-07-15 03:07:12}
\Level{A}
\SubArea{Elementary arithmetics}
\LANGen{
\Question{
The result of subtraction \( 0.36-3.21 \) is:}
{\( -2.85 \)}{\( 2.85 \)}{\( 3.57 \)}{\( 3.15 \)}{}{}}
\LANGpl{
\Question{
Wartością wyrażenia \( 0{,}36-3{,}21 \) jest:}
{\( -2{,}85 \)}{\( 2{,}85 \)}{\( 3{,}57 \)}{\( 3{,}15 \)}{}{}}
\LANGcs{
\Question{
Hodnota výrazu \( 0{,}36-3{,}21 \) je:}
{\( -2{,}85 \)}{\( 2{,}85 \)}{\( 3{,}57 \)}{\( 3{,}15 \)}{}{}}
\LANGsk{
\Question{
Hodnota výrazu \( 0{,}36-3{,}21 \) je:}
{\( -2{,}85 \)}{\( 2{,}85 \)}{\( 3{,}57 \)}{\( 3{,}15 \)}{}{}}
\LANGes{
\Question{
El resultado de la resta \( 0.36-3.21 \) es:}
{\( -2.85 \)}{\( 2.85 \)}{\( 3.57 \)}{\( 3.15 \)}{}{}}
\End

\Data\ID{1003034105}
\Updated{2020-07-15 03:07:12}
\Level{A}
\SubArea{Elementary arithmetics}
\LANGen{
\Question{
Evaluate the expression \( 0.864:0.24 \).}
{\( 3.6 \)}{\( 36 \)}{\( 0.36 \)}{\( 0.624 \)}{}{}}
\LANGpl{
\Question{
Wartością wyrażenia \( 0{,}864:0{,}24 \) jest:}
{\( 3{,}6 \)}{\( 36 \)}{\( 0{,}36 \)}{\( 0{,}624 \)}{}{}}
\LANGcs{
\Question{
Určete hodnotu výrazu \( 0{,}864:0{,}24 \).}
{\( 3{,}6 \)}{\( 36 \)}{\( 0{,}36 \)}{\( 0{,}624 \)}{}{}}
\LANGsk{
\Question{
Určte hodnotu výrazu \( 0{,}864:0{,}24 \).}
{\( 3{,}6 \)}{\( 36 \)}{\( 0{,}36 \)}{\( 0{,}624 \)}{}{}}
\LANGes{
\Question{
Calcula la expresión  \( 0.864:0.24 \).}
{\( 3.6 \)}{\( 36 \)}{\( 0.36 \)}{\( 0.624 \)}{}{}}
\End

\Data\ID{1003034109}
\Updated{2020-07-15 03:07:12}
\Level{A}
\SubArea{Elementary arithmetics}
\LANGen{
\Question{
Give the result of   \( 3+5-2\cdot3+\left(2^3-6\right):1\).}
{\( 4 \)}{\( 20 \)}{\( 14 \)}{\( 2 \)}{}{}}
\LANGpl{
\Question{
Wartość wyrażenia  \( 3+5-2\cdot3+\left(2^3-6\right):1\) jest równa:}
{\( 4 \)}{\( 20 \)}{\( 14 \)}{\( 2 \)}{}{}}
\LANGcs{
\Question{
Vypočítejte \( 3+5-2\cdot3+\left(2^3-6\right):1\).}
{\( 4 \)}{\( 20 \)}{\( 14 \)}{\( 2 \)}{}{}}
\LANGsk{
\Question{
Vypočítajte \( 3+5-2\cdot3+\left(2^3-6\right):1\).}
{\( 4 \)}{\( 20 \)}{\( 14 \)}{\( 2 \)}{}{}}
\LANGes{
\Question{
Calcula  \( 3+5-2\cdot3+\left(2^3-6\right):1\).}
{\( 4 \)}{\( 20 \)}{\( 14 \)}{\( 2 \)}{}{}}
\End

\Data\ID{1003034110}
\Updated{2020-07-15 03:07:12}
\Level{A}
\SubArea{Elementary arithmetics}
\LANGen{
\Question{
Calculating \(2\cdot3-(4\cdot15-36)\) we get:}
{\( -18 \)}{\( 18 \)}{\( 90 \)}{\( -90 \)}{}{}}
\LANGpl{
\Question{
Wartością wyrażenia \(2\cdot3-(4\cdot15-36)\) jest:}
{\( -18 \)}{\( 18 \)}{\( 90 \)}{\( -90 \)}{}{}}
\LANGcs{
\Question{
Vypočítejte \(2\cdot3-(4\cdot15-36)\).}
{\( -18 \)}{\( 18 \)}{\( 90 \)}{\( -90 \)}{}{}}
\LANGsk{
\Question{
Vypočítajte \(2\cdot3-(4\cdot15-36)\).}
{\( -18 \)}{\( 18 \)}{\( 90 \)}{\( -90 \)}{}{}}
\LANGes{
\Question{
Al calcular  \(2\cdot3-(4\cdot15-36)\) obtenemos:}
{\( -18 \)}{\( 18 \)}{\( 90 \)}{\( -90 \)}{}{}}
\End

\Data\ID{1003099201}
\Updated{2021-07-24 10:07:47}
\Level{A}
\SubArea{Elementary arithmetics}
\LANGen{
\Question{
A toe sometimes appears as a unit of energy in various publications. One toe is a unit of energy defined as the amount of energy released by burning one tonne of crude oil and it equals \( 41 868\,\mathrm{MJ} \) (\( 1 \,\mathrm{MJ} = 1\, 000\, 000 \,\mathrm{J} \)). How many joules are in \( 1 \) toe?}
{\( 4.1868\cdot10^{10} \)}{\( 4.1868\cdot10^{8} \)}{\( 4.1868\cdot10^{9} \)}{\( 4.1868\cdot10^{11} \)}{}{}}
\LANGpl{
\Question{
W różnych publikacjach jako jednostka energii pojawia się czasem toe. \(1\) toe odpowiada energii jaką uzyskuje się z \( 1\) tony ropy naftowej i równa się \( 41 868\,\mathrm{MJ} \) (\( 1 \,\mathrm{MJ} = 1\, 000\, 000 \,\mathrm{J} \)). Ilu dżulom równa się  \( 1 \) toe?}
{\( 4{,}1868\cdot10^{10} \)}{\( 4{,}1868\cdot10^{8} \)}{\( 4{,}1868\cdot10^{9} \)}{\( 4{,}1868\cdot10^{11} \)}{}{}}
\LANGcs{
\Question{
V publikacích se jako jednotka energie někdy objevuje tzv. tuna olejového ekvivalentu (toe). Jeden toe je jednotka energie definovaná jako množství energie uvolněné spálením jedné tuny ropy a rovná se \( 41\, 868\,\mathrm{MJ} \) (\( 1 \,\mathrm{MJ} = 1\, 000\, 000 \,\mathrm{J} \)). Kolik  joulů je \( 1 \) toe?}
{\( 4{,}1868\cdot10^{10} \)}{\( 4{,}1868\cdot10^{8} \)}{\( 4{,}1868\cdot10^{9} \)}{\( 4{,}1868\cdot10^{11} \)}{}{}}
\LANGsk{
\Question{
V rôznych publikáciach sa ako jednotka energie niekedy objavuje tzv. tona olejového ekvivalentu (toe). Jeden toe sa ako jednotka energie definuje ako množstvo energie uvolnené spálením jednej tony ropy a rovná sa \( 41\, 868\,\mathrm{MJ} \) (\( 1 \,\mathrm{MJ} = 1\, 000\, 000 \,\mathrm{J} \)). Koľko joulov má \( 1 \) toe?}
{\( 4{,}1868\cdot10^{10} \)}{\( 4{,}1868\cdot10^{8} \)}{\( 4{,}1868\cdot10^{9} \)}{\( 4{,}1868\cdot10^{11} \)}{}{}}
\LANGes{
\Question{
En varias publicaciones aparece como unidad de energía la tonelada equivalente de petróleo (toe). Un toe es una unidad de energía definida como la cantidad de energía liberada al quemar una tonelada de petróleo crudo y equivale a \( 41 868\,\mathrm{MJ} \) (\( 1 \,\mathrm{MJ} = 1\, 000\, 000 \,\mathrm{J} \)). ¿Cuántos julios hay en \( 1 \) toe?}
{\( 4.1868\cdot10^{10} \)}{\( 4.1868\cdot10^{8} \)}{\( 4.1868\cdot10^{9} \)}{\( 4.1868\cdot10^{11} \)}{}{}}
\End

\Data\ID{1003099202}
\Updated{2020-08-28 02:08:00}
\Level{A}
\SubArea{Elementary arithmetics}
\LANGen{
\Question{
The average distance from Mars to the Sun is \( 2.28\cdot10^8\,\mathrm{km} \).  Give the distance in the standard form.}
{\( 228\,000\,000\,\mathrm{km} \)}{\( 22\,800\,000\,\mathrm{km} \)}{\( 2\,280\,000\,000\,\mathrm{km} \)}{\( 22\,800\,000\,000\,\mathrm{km} \)}{}{}}
\LANGpl{
\Question{
Średnia odległość Marsa od Słońca wynosi \( 2{,}28\cdot10^8\,\mathrm{km} \). Odległość ta zapisana bez użycia potęgi jest równa:}
{\( 228\, 000\, 000\,\mathrm{km} \)}{\( 22\, 800\, 000\,\mathrm{km} \)}{\( 2\, 280\, 000\, 000\,\mathrm{km} \)}{\( 22\, 800\, 000\, 000\,\mathrm{km} \)}{}{}}
\LANGcs{
\Question{
Průměrná vzdálenost mezi Marsem a Zemí je \( 2{,}28\cdot10^8\,\mathrm{km} \).  Uveďte vzdálenost ve standardizované formě zápisu čísla.}
{\( 228\, 000\, 000\,\mathrm{km} \)}{\( 22\, 800\, 000\,\mathrm{km} \)}{\( 2\, 280\, 000\, 000\,\mathrm{km} \)}{\( 22\, 800\, 000\, 000\,\mathrm{km} \)}{}{}}
\LANGsk{
\Question{
Priemerná vzdialenosť medzi Marsom a Zemou je \( 2{,}28\cdot10^8\,\mathrm{km} \). Uveďte vzdialenosť v štandardnej forme zápisu čísla.}
{\( 228\, 000\, 000\,\mathrm{km} \)}{\( 22\, 800\, 000\,\mathrm{km} \)}{\( 2\, 280\, 000\, 000\,\mathrm{km} \)}{\( 22\, 800\, 000\, 000\,\mathrm{km} \)}{}{}}
\LANGes{
\Question{
La distancia media entre Marte y el Sol es de \( 2.28\cdot10^8\,\mathrm{km} \). Indica la distancia en forma estándar.}
{\( 228\,000\,000\,\mathrm{km} \)}{\( 22\,800\,000\,\mathrm{km} \)}{\( 2\,280\,000\,000\,\mathrm{km} \)}{\( 22\,800\,000\,000\,\mathrm{km} \)}{}{}}
\End

\Data\ID{1003099203}
\Updated{2020-07-15 03:07:12}
\Level{A}
\SubArea{Elementary arithmetics}
\LANGen{
\Question{
A man sleeps \( 8 \) hours a day on average. How many seconds on average does he sleep during the longest month of the year? Give the result in exponential notation.}
{\( 8.928\cdot10^5\,\mathrm{s} \)}{\( 1.488\cdot10^4\,\mathrm{s} \)}{\( 8928\cdot10^2\,\mathrm{s} \)}{\( 8.64\cdot10^5\,\mathrm{s} \)}{}{}}
\LANGpl{
\Question{
Człowiek śpi średnio  \( 8 \) godzin w ciągu doby. Ile sekund przesypia w ciągu najdłuższego miesiąca? Wynik podaj w notacji wykładniczej.}
{\( 8{,}928\cdot10^5\,\mathrm{s} \)}{\( 1{,}488\cdot10^4\,\mathrm{s} \)}{\( 8928\cdot10^2\,\mathrm{s} \)}{\( 8{,}64\cdot10^5\,\mathrm{s} \)}{}{}}
\LANGcs{
\Question{
Člověk spí průměrně \( 8 \) hodin denně. Průměrně kolik sekund spí člověk po dobu nejdelšího měsíce roku? Výsledek uveďte ve vědeckém tvaru.}
{\( 8{,}928\cdot10^5\,\mathrm{s} \)}{\( 1{,}488\cdot10^4\,\mathrm{s} \)}{\( 8928\cdot10^2\,\mathrm{s} \)}{\( 8{,}64\cdot10^5\,\mathrm{s} \)}{}{}}
\LANGsk{
\Question{
Človek spí priemerne \( 8 \) hodín denne. Priemerne koľko sekúnd spí človek počas najdlhšieho mesiaca v roku? Výsledok uveďte vo vedeckom tvare.}
{\( 8{,}928\cdot10^5\,\mathrm{s} \)}{\( 1{,}488\cdot10^4\,\mathrm{s} \)}{\( 8928\cdot10^2\,\mathrm{s} \)}{\( 8{,}64\cdot10^5\,\mathrm{s} \)}{}{}}
\LANGes{
\Question{
Un hombre duerme  \( 8 \) horas al día de media. ¿Cuántos segundos duerme en promedio durante el mes más largo del año? Indica el resultado en notación científica.}
{\( 8.928\cdot10^5\,\mathrm{s} \)}{\( 1.488\cdot10^4\,\mathrm{s} \)}{\( 8928\cdot10^2\,\mathrm{s} \)}{\( 8.64\cdot10^5\,\mathrm{s} \)}{}{}}
\End

\Data\ID{1003099204}
\Updated{2020-07-15 03:07:12}
\Level{A}
\SubArea{Elementary arithmetics}
\LANGen{
\Question{
Consider the information given in the table. Calculate the shortest distance between Neptune (the farthest planet from the Sun) and Mercury (the closest planet to the Sun).
\[
\begin{array}{|c|c|} \hline \text{Planet} & \text{Distance from the Sun } (\mathrm{km} ) \\\hline \text{Mercury} & 5.8\cdot10^7 \\\hline \text{Neptune} & 4.5\cdot10^9 \\\hline \end{array}
\]}
{\( 4.442\cdot10^9\,\mathrm{km} \)}{\( 4.558\cdot10^9\,\mathrm{km} \)}{\( 1.3\cdot10^7\,\mathrm{km} \)}{\( 1.3\cdot10^9\,\mathrm{km} \)}{}{}}
\LANGpl{
\Question{
Korzystając z informacji podanych w tabeli oblicz najkrótszą odległość między Neptunem i Merkurym -  najbardziej oddaloną planetą od Słońca i planetą leżąca najbliżej Słońca.  \[\begin{array}{|c|c|} \hline \text{Planeta} & \text{Odległość od Słońca } (\mathrm{km} ) \\\hline \text{Merkury} & 5{,}8\cdot10^7 \\\hline \text{Neptun} & 4{,}5\cdot10^9 \\\hline \end{array}
\]}
{\( 4{,}442\cdot10^9\,\mathrm{km} \)}{\( 4{,}558\cdot10^9\,\mathrm{km} \)}{\( 1{,}3\cdot10^7\,\mathrm{km} \)}{\( 1{,}3\cdot10^9\,\mathrm{km} \)}{}{}}
\LANGcs{
\Question{
Použijte informace uvedené v tabulce. Vypočítejte nejkratší vzdálenost mezi Neptunem (nejvzdálenější planetou od Slunce) a Merkurem (nejbližší planetou k Slunci).
\[
\begin{array}{|c|c|} \hline \text{Planeta} & \text{Vzdálenost od Slunce  } (\mathrm{km} ) \\\hline \text{Merkur} & 5{,}8\cdot10^7 \\\hline \text{Neptun} & 4{,}5\cdot10^9 \\\hline \end{array}
\]}
{\( 4{,}442\cdot10^9\,\mathrm{km} \)}{\( 4{,}558\cdot10^9\,\mathrm{km} \)}{\( 1{,}3\cdot10^7\,\mathrm{km} \)}{\( 1{,}3\cdot10^9\,\mathrm{km} \)}{}{}}
\LANGsk{
\Question{
Použite informácie uvedené v tabuľke. Vypočítajte najkratšiu vzdialenosť medzi Neptúnom (najvzdialenejšou planétou od Slnka) a Merkúrom (najbližšou planétou k Slnku).
\[
\begin{array}{|c|c|} \hline \text{Planéta} & \text{Vzdialenosť od Slnka } (\mathrm{km} ) \\\hline \text{Merkúr} & 5{,}8\cdot10^7 \\\hline \text{Neptún} & 4{,}5\cdot10^9 \\\hline \end{array}
\]}
{\( 4{,}442\cdot10^9\,\mathrm{km} \)}{\( 4{,}558\cdot10^9\,\mathrm{km} \)}{\( 1{,}3\cdot10^7\,\mathrm{km} \)}{\( 1{,}3\cdot10^9\,\mathrm{km} \)}{}{}}
\LANGes{
\Question{
Considerando la información dada en la tabla, calcula la distancia más corta entre Neptuno (el planeta más lejano del Sol) y Mercurio (el planeta más cercano al Sol).
\[
\begin{array}{|c|c|} \hline \text{Planet} & \text{Distance from the Sun } (\mathrm{km} ) \\\hline \text{Mercury} & 5.8\cdot10^7 \\\hline \text{Neptune} & 4.5\cdot10^9 \\\hline \end{array}
\]}
{\( 4.442\cdot10^9\,\mathrm{km} \)}{\( 4.558\cdot10^9\,\mathrm{km} \)}{\( 1.3\cdot10^7\,\mathrm{km} \)}{\( 1.3\cdot10^9\,\mathrm{km} \)}{}{}}
\End

\Data\ID{1003099205}
\Updated{2020-07-15 03:07:12}
\Level{A}
\SubArea{Elementary arithmetics}
\LANGen{
\Question{
Mass of a hydrogen atom is  \( 0.000\,000\,000\,000\,000\,000\,000\,001\,67\,\mathrm{g} \). Give the mass in exponential notation.}
{\( 1.67\cdot10^{-24}\,\mathrm{g} \)}{\( 167\cdot10^{-26}\,\mathrm{g} \)}{\( 1.67\cdot10^{24}\,\mathrm{g} \)}{\( 167\cdot10^{-21}\,\mathrm{g} \)}{}{}}
\LANGpl{
\Question{
Masa atomu wodoru wynosi  \( 0{,}000\,000\,000\,000\,000\,000\,000\,001\,67\,\mathrm{g} \). Podaj masę w notacji wykładniczej.}
{\( 1{,}67\cdot10^{-24}\,\mathrm{g} \)}{\( 167\cdot10^{-26}\,\mathrm{g} \)}{\( 1{,}67\cdot10^{24}\,\mathrm{g} \)}{\( 167\cdot10^{-21}\,\mathrm{g} \)}{}{}}
\LANGcs{
\Question{
Atomová hmotnost vodíku je \( 0{,}000\,000\,000\,000\,000\,000\,000\,001\,67\,\mathrm{g} \). Uveďte danou hmotnost pomocí vědeckého zápisu čísel.}
{\( 1{,}67\cdot10^{-24}\,\mathrm{g} \)}{\( 167\cdot10^{-26}\,\mathrm{g} \)}{\( 1{,}67\cdot10^{24}\,\mathrm{g} \)}{\( 167\cdot10^{-21}\,\mathrm{g} \)}{}{}}
\LANGsk{
\Question{
Atómová hmotnosť vodíka je \( 0{,}000\,000\,000\,000\,000\,000\,000\,001\,67\,\mathrm{g} \). Uveďte danú hmotnosť pomocou vedeckého zápisu čísel.}
{\( 1{,}67\cdot10^{-24}\,\mathrm{g} \)}{\( 167\cdot10^{-26}\,\mathrm{g} \)}{\( 1{,}67\cdot10^{24}\,\mathrm{g} \)}{\( 167\cdot10^{-21}\,\mathrm{g} \)}{}{}}
\LANGes{
\Question{
La masa atómica del hidrógeno es de  \( 0.000\,000\,000\,000\,000\,000\,000\,001\,67\,\mathrm{g} \). Indica la masa atómica en notación científica.}
{\( 1.67\cdot10^{-24}\,\mathrm{g} \)}{\( 167\cdot10^{-26}\,\mathrm{g} \)}{\( 1.67\cdot10^{24}\,\mathrm{g} \)}{\( 167\cdot10^{-21}\,\mathrm{g} \)}{}{}}
\End

\Data\ID{1003099206}
\Updated{2020-07-15 03:07:12}
\Level{A}
\SubArea{Elementary arithmetics}
\LANGen{
\Question{
Earth weighs \( 5.98\cdot10^{24}\,\mathrm{kg} \), the Moon weighs \( 7\cdot10^{22}\,\mathrm{kg} \). Calculate in kilograms the weight of both the Moon and Earth.}
{\( 6.05\cdot10^{24} \)}{\( 6.05\cdot10^{22} \)}{\( 1.298\cdot10^{24} \)}{\( 12.98\cdot10^{24} \)}{}{}}
\LANGpl{
\Question{
Ziemia waży \( 5{,}98\cdot10^{24}\,\mathrm{kg} \), Księżyc waży \( 7\cdot10^{22}\,\mathrm{kg} \). Ile kilogramów ważą w sumie Ziemia i Księżyc?}
{\( 6{,}05\cdot10^{24} \)}{\( 6{,}05\cdot10^{22} \)}{\( 1{,}298\cdot10^{24} \)}{\( 12{,}98\cdot10^{24} \)}{}{}}
\LANGcs{
\Question{
Hmotnost Země je \( 5{,}98\cdot10^{24}\,\mathrm{kg} \) a hmotnost Měsíce je \( 7\cdot10^{22}\,\mathrm{kg} \). Určete (v kilogramech) součet jejich hmotností a zapište výsledek ve vědeckém zápise.}
{\( 6{,}05\cdot10^{24} \)}{\( 6{,}05\cdot10^{22} \)}{\( 1{,}298\cdot10^{24} \)}{\( 12{,}98\cdot10^{24} \)}{}{}}
\LANGsk{
\Question{
Hmotnosť Zeme je \( 5{,}98\cdot10^{24}\,\mathrm{kg} \) a hmotnosť Mesiaca je \( 7\cdot10^{22}\,\mathrm{kg} \). Určte (v kilogramoch) súčet ich hmotností a zapíšte výsledok vo vedeckom zápise.}
{\( 6{,}05\cdot10^{24} \)}{\( 6{,}05\cdot10^{22} \)}{\( 1{,}298\cdot10^{24} \)}{\( 12{,}98\cdot10^{24} \)}{}{}}
\LANGes{
\Question{
La Tierra pesa  \( 5.98\cdot10^{24}\,\mathrm{kg} \), la Luna pesa  \( 7\cdot10^{22}\,\mathrm{kg} \). Calcula el peso de la Luna y la Tierra en kilogramos.}
{\( 6.05\cdot10^{24} \)}{\( 6.05\cdot10^{22} \)}{\( 1.298\cdot10^{24} \)}{\( 12.98\cdot10^{24} \)}{}{}}
\End

\Data\ID{1003118609}
\Updated{2020-07-15 03:07:12}
\Level{A}
\SubArea{Elementary arithmetics}
\LANGen{
\Question{
The number \( 350 \) millions is equal to:}
{\( 3.5\cdot10^8 \)}{\( 3.5\cdot10^6 \)}{\( 3.5\cdot10^9 \)}{\( 3.5\cdot10^7 \)}{}{}}
\LANGpl{
\Question{
Liczba \( 350 \) milionów jest równa:}
{\( 3{,}5\cdot10^8 \)}{\( 3{,}5\cdot10^6 \)}{\( 3{,}5\cdot10^9 \)}{\( 3{,}5\cdot10^7 \)}{}{}}
\LANGcs{
\Question{
Číslo \( 350 \) milionů je rovno:}
{\( 3{,}5\cdot10^8 \)}{\( 3{,}5\cdot10^6 \)}{\( 3{,}5\cdot10^9 \)}{\( 3{,}5\cdot10^7 \)}{}{}}
\LANGsk{
\Question{
Číslo \( 350 \) miliónov sa rovná:}
{\( 3{,}5\cdot10^8 \)}{\( 3{,}5\cdot10^6 \)}{\( 3{,}5\cdot10^9 \)}{\( 3{,}5\cdot10^7 \)}{}{}}
\LANGes{
\Question{
El número \( 350 \) millones equivale a:}
{\( 3.5\cdot10^8 \)}{\( 3.5\cdot10^6 \)}{\( 3.5\cdot10^9 \)}{\( 3.5\cdot10^7 \)}{}{}}
\End

\Data\ID{1003118610}
\Updated{2020-07-15 03:07:12}
\Level{A}
\SubArea{Elementary arithmetics}
\LANGen{
\Question{
Which of the following numbers is the smallest?}
{\( 1.23\cdot10^{-5} \)}{\( 25.36\cdot10^{-4} \)}{\( 0.00000013\cdot10^5 \)}{\( \frac{21}{0.7\cdot10^5} \)}{}{}}
\LANGpl{
\Question{
Która z podanych liczb jest najmniejsza?}
{\( 1{,}23\cdot10^{-5} \)}{\( 25{,}36\cdot10^{-4} \)}{\( 0{,}00000013\cdot10^5 \)}{\( \frac{21}{0{,}7\cdot10^5} \)}{}{}}
\LANGcs{
\Question{
Které z uvedených čísel je nejmenší?}
{\( 1{,}23\cdot10^{-5} \)}{\( 25{,}36\cdot10^{-4} \)}{\( 0{,}00000013\cdot10^5 \)}{\( \frac{21}{0{,}7\cdot10^5} \)}{}{}}
\LANGsk{
\Question{
Ktoré z nasledujúcich čísel je najmenšie?}
{\( 1{,}23\cdot10^{-5} \)}{\( 25{,}36\cdot10^{-4} \)}{\( 0{,}00000013\cdot10^5 \)}{\( \frac{21}{0{,}7\cdot10^5} \)}{}{}}
\LANGes{
\Question{
¿Cuál de los números siguientes es el más pequeño?}
{\( 1.23\cdot10^{-5} \)}{\( 25.36\cdot10^{-4} \)}{\( 0.00000013\cdot10^5 \)}{\( \frac{21}{0.7\cdot10^5} \)}{}{}}
\End

\Data\ID{9000085601}
\Updated{2020-07-15 03:07:37}
\Level{A}
\SubArea{Elementary arithmetics}
\LANGen{
\Question{
Evaluate the following numbers and round to the nearest tens. 
\[ 
 10^{2} + 11^{2} + 12^{2} + 13^{2} + 14^{2}
\]}
{\(730\)}{\(720\)}{\(740\)}{\(750\)}{}{}}
\LANGpl{
\Question{
Liczba  
\[ 
 10^{2} + 11^{2} + 12^{2} + 13^{2} + 14^{2}
\] po zaokrągleniu do dziesiątek jest równa:}
{\(730\)}{\(720\)}{\(740\)}{\(750\)}{}{}}
\LANGcs{
\Question{
Číslo \(10^{2} + 11^{2} + 12^{2} + 13^{2} + 14^{2}\) 
je po zaokrouhlení na desítky rovno:}
{\(730\)}{\(720\)}{\(740\)}{\(750\)}{}{}}
\LANGsk{
\Question{
Číslo \(10^{2} + 11^{2} + 12^{2} + 13^{2} + 14^{2}\) 
sa po zaokrúhlení na desiatky rovná:}
{\(730\)}{\(720\)}{\(740\)}{\(750\)}{}{}}
\LANGes{
\Question{
Evalúa los números siguientes y redondea a las decenas más cercanas.
\[ 
 10^{2} + 11^{2} + 12^{2} + 13^{2} + 14^{2}
\]}
{\(730\)}{\(720\)}{\(740\)}{\(750\)}{}{}}
\End

\Data\ID{9000085603}
\Updated{2020-07-15 03:07:37}
\Level{A}
\SubArea{Elementary arithmetics}
\LANGen{
\Question{
Find the sum of the three numbers obtained by rounding the number
\(5\: 316\) to the
nearest tens, hundreds and thousands.}
{\(15\: 620\)}{\(15\: 610\)}{\(15\: 560\)}{\(15\: 580\)}{}{}}
\LANGpl{
\Question{
Znajdź sumę trzech liczb uzyskanych przez zaokrąglenie liczby
\(5\: 316\) do dziesiątek, setek i tysięcy.}
{\(15\: 620\)}{\(15\: 610\)}{\(15\: 560\)}{\(15\: 580\)}{}{}}
\LANGcs{
\Question{
Určete součet tří čísel, která získáme, zaokrouhlíme-li číslo
\(5\: 316\) na
desítky, na stovky a na tisíce.}
{\(15\: 620\)}{\(15\: 610\)}{\(15\: 560\)}{\(15\: 580\)}{}{}}
\LANGsk{
\Question{
Určte súčet troch čísel, ktoré dostaneme zaokrúhlením čísla 
\(5\: 316\) na desiatky, stovky a tisícky.}
{\(15\: 620\)}{\(15\: 610\)}{\(15\: 560\)}{\(15\: 580\)}{}{}}
\LANGes{
\Question{
Calcula la suma de los tres números obtenidos al redondear el número 
\(5\: 316\) a las decenas, centenas y millares más cercanos.}
{\(15\: 620\)}{\(15\: 610\)}{\(15\: 560\)}{\(15\: 580\)}{}{}}
\End

\Data\ID{9000085604}
\Updated{2020-07-15 03:07:37}
\Level{A}
\SubArea{Elementary arithmetics}
\LANGen{
\Question{
Find the sum of the three numbers obtained by rounding the number
\(2\: 013\) to the
nearest tens, hundreds and thousands.}
{\(6\: 010\)}{\(6\: 000\)}{\(6\: 020\)}{\(6\: 030\)}{}{}}
\LANGpl{
\Question{
Znajdź sumę trzech liczb uzyskanych przez zaokrąglenie liczby
\(2\: 013\) do dziesiątek, setek i tysięcy.}
{\(6\: 010\)}{\(6\: 000\)}{\(6\: 020\)}{\(6\: 030\)}{}{}}
\LANGcs{
\Question{
Určete součet tří čísel, která získáme, zaokrouhlíme-li číslo
\(2\: 013\) na
desítky, na stovky a na tisíce.}
{\(6\: 010\)}{\(6\: 000\)}{\(6\: 020\)}{\(6\: 030\)}{}{}}
\LANGsk{
\Question{
Určte súčet troch čísel, ktoré dostaneme zaokrúhlením čísla \(2\: 013\)  na desiatky, stovky a tisícky.}
{\(6\: 010\)}{\(6\: 000\)}{\(6\: 020\)}{\(6\: 030\)}{}{}}
\LANGes{
\Question{
Calcula la suma de los tres números obtenidos al redondear el número 
\(2\: 013\) a las decenas, centenas y millares más cercanos.}
{\(6\: 010\)}{\(6\: 000\)}{\(6\: 020\)}{\(6\: 030\)}{}{}}
\End

\Data\ID{9000085607}
\Updated{2020-07-15 03:07:37}
\Level{A}
\SubArea{Elementary arithmetics}
\LANGen{
\Question{
Given numbers \(8\: 175\) 
and \(3\: 926\),
find the difference between  the sum of these numbers rounded to the nearest tens and the sum of these numbers rounded to the nearest hundreds.}
{\(10\)}{\(20\)}{\(100\)}{\(0\)}{}{}}
\LANGpl{
\Question{
Dla podanych liczb \(8\: 175\) 
i  \(3\: 926\),
znajdź różnicę pomiędzy sumą tych liczb zaokrągloną do dziesiątek i sumą tych liczb zaokrągloną do setek.}
{\(10\)}{\(20\)}{\(100\)}{\(0\)}{}{}}
\LANGcs{
\Question{
Jsou dána čísla \(8\: 175\) 
a \(3\: 926\). O
kolik je součet těchto dvou čísel zaokrouhlených na desítky větší než
součet těchto dvou čísel zaokrouhlených na stovky?}
{\(10\)}{\(20\)}{\(100\)}{\(0\)}{}{}}
\LANGsk{
\Question{
Dané sú čísla \(8\: 175\) 
a \(3\: 926\). Určte rozdiel medzi súčtom daných čísel zaokrúhlených na desiatky a súčtom daných čísel zaokrúhlených na stovky.}
{\(10\)}{\(20\)}{\(100\)}{\(0\)}{}{}}
\LANGes{
\Question{
Dados los números \(8\: 175\) 
y \(3\: 926\),
calcula la diferencia entre la suma de estos números redondeados a las decenas más cercanas y la suma de estos números redondeados a las centenas más cercanas.}
{\(10\)}{\(20\)}{\(100\)}{\(0\)}{}{}}
\End

\Data\ID{9000085608}
\Updated{2020-07-15 03:07:37}
\Level{A}
\SubArea{Elementary arithmetics}
\LANGen{
\Question{
Given numbers \(456\: 138\) 
and \(321\: 814\),
find the difference between the sum of these numbers rounded to the nearest tens and the sum of these numbers rounded to the nearest hundreds.}
{\(50\)}{\(100\)}{\(1\: 000\)}{\(0\)}{}{}}
\LANGpl{
\Question{
Dla danych liczb \(456\: 138\) 
i \(321\: 814\),
znajdź różnicę pomiędzy sumą tych liczb zaokrągloną do dziesiątek i  sumą tych liczb zaokrągloną do setek.}
{\(50\)}{\(100\)}{\(1\: 000\)}{\(0\)}{}{}}
\LANGcs{
\Question{
Jsou dána čísla \(456\: 138\) 
a \(321\: 814\). O
kolik je součet těchto dvou čísel zaokrouhlených na desítky větší než
součet těchto dvou čísel zaokrouhlených na stovky?}
{\(50\)}{\(100\)}{\(1\: 000\)}{\(0\)}{}{}}
\LANGsk{
\Question{
Dané sú čísla \(456\: 138\) 
a \(321\: 814\). Určte rozdiel medzi súčtom daných čísel zaokrúhlených na desiatky a súčtom daných čísel zaokrúhlených na stovky.}
{\(50\)}{\(100\)}{\(1\: 000\)}{\(0\)}{}{}}
\LANGes{
\Question{
Dados los números  \(456\: 138\) 
y \(321\: 814\),
calcula la diferencia entre la suma de estos números redondeados a las decenas más cercanas y la suma de estos números redondeados a las centenas más cercanas.}
{\(50\)}{\(100\)}{\(1\: 000\)}{\(0\)}{}{}}
\End

\Data\ID{9000085609}
\Updated{2020-07-15 03:07:37}
\Level{A}
\SubArea{Elementary arithmetics}
\LANGen{
\Question{
Given the number \(45\: 875\),
round this number to nearest thousands, nearest hundreds and subtract the results.}
{\(100\)}{\(200\)}{\(1\: 000\)}{\(0\)}{}{}}
\LANGpl{
\Question{
Daną liczbę  \(45\: 875\),
zaokrągli do tysiąca oraz do setek i oblicz różnicę tych liczb.}
{\(100\)}{\(200\)}{\(1\: 000\)}{\(0\)}{}{}}
\LANGcs{
\Question{
Je dáno číslo \(45\: 875\).
O kolik bude toto číslo zaokrouhlené na tisíce větší než toto číslo
zaokrouhlené na stovky?}
{\(100\)}{\(200\)}{\(1\: 000\)}{\(0\)}{}{}}
\LANGsk{
\Question{
Dané je číslo \(45\: 875\). O koľko bude toto číslo zaokrúhlené na tisícky väčšie ako toto číslo zaokrúhlené na stovky?}
{\(100\)}{\(200\)}{\(1\: 000\)}{\(0\)}{}{}}
\LANGes{
\Question{
Dado el número \(45\: 875\),
redondea este número a los millares más cercanos, a las centenas más cercanas y resta los resultados.}
{\(100\)}{\(200\)}{\(1\: 000\)}{\(0\)}{}{}}
\End

\Data\ID{9000085610}
\Updated{2020-07-15 03:07:37}
\Level{A}
\SubArea{Elementary arithmetics}
\LANGen{
\Question{
Given the number \(82\: 361\),
round this number to nearest thousands, nearest hundreds and subtract the results.}
{\(400\)}{\(300\)}{\(200\)}{\(100\)}{}{}}
\LANGpl{
\Question{
Daną liczbę \(82\: 361\),
zaokrągli do tysiąca oraz do setek i oblicz różnicę tych liczb.}
{\(400\)}{\(300\)}{\(200\)}{\(100\)}{}{}}
\LANGcs{
\Question{
Je dáno číslo \(82\: 361\).
O kolik bude toto číslo zaokrouhlené na tisíce menší než toto číslo
zaokrouhlené na stovky?}
{\(400\)}{\(300\)}{\(200\)}{\(100\)}{}{}}
\LANGsk{
\Question{
Dané je číslo \(82\: 361\). O koľko bude toto číslo zaokrúhlené na tisícky menšie ako toto číslo zaokrúhlené na stovky?}
{\(400\)}{\(300\)}{\(200\)}{\(100\)}{}{}}
\LANGes{
\Question{
Dado el número \(82\: 361\),
redondea este número a los millares más cercanos, a las centenas más cercanas y resta los resultados.}
{\(400\)}{\(300\)}{\(200\)}{\(100\)}{}{}}
\End

\Data\ID{1003099401}
\Updated{2020-07-15 03:07:21}
\Level{B}
\SubArea{Elementary arithmetics}
\LANGen{
\Question{
The number \( 43256232a2 \) is divisible by \( 9 \) if}
{\( a= 7 \).}{\( a= 1 \).}{\( a= 0 \).}{\( a= 4 \).}{}{}}
\LANGpl{
\Question{
Liczba \( 43256232a2 \) jest podzielna przez \( 9 \) jeżeli:}
{\( a= 7 \).}{\( a= 1 \).}{\( a= 0 \).}{\( a= 4 \).}{}{}}
\LANGcs{
\Question{
Číslo \( 43256232a2 \) je dělitelné \( 9 \), jestliže}
{\( a= 7 \).}{\( a= 1 \).}{\( a= 0 \).}{\( a= 4 \).}{}{}}
\LANGsk{
\Question{
Číslo \( 43256232a2 \) je deliteľné \( 9 \), ak}
{\( a= 7 \).}{\( a= 1 \).}{\( a= 0 \).}{\( a= 4 \).}{}{}}
\LANGes{
\Question{
El número  \( 43256232a2 \) es divisible por  \( 9 \) si}
{\( a= 7 \).}{\( a= 1 \).}{\( a= 0 \).}{\( a= 4 \).}{}{}}
\End

\Data\ID{1003099402}
\Updated{2020-07-15 03:07:21}
\Level{B}
\SubArea{Elementary arithmetics}
\LANGen{
\Question{
How many odd double-digit numbers give the reminder of \( 2 \) when divided by \( 9 \) and are divisible by \( 13 \) as well?}
{\( 1 \)}{\( 2 \)}{\( 3 \)}{\( 4 \)}{}{}}
\LANGpl{
\Question{
Ile jest liczb dwucyfrowych nieparzystych, które przy dzieleniu przez \( 9 \) dają resztę \( 2 \) i jednocześnie są podzielne przez  \( 13 \)?}
{\( 1 \)}{\( 2 \)}{\( 3 \)}{\( 4 \)}{}{}}
\LANGcs{
\Question{
Kolik sudých dvouciferných čísel má zbytek \( 2 \), když je vydělíme \( 9 \) a zároveň jsou dělitelné \( 13 \)?}
{\( 1 \)}{\( 2 \)}{\( 3 \)}{\( 4 \)}{}{}}
\LANGsk{
\Question{
Koľko párnych dvojciferných čísel po vydelení \( 9 \) dáva zvyšok \( 2 \) a zároveň je deliteľné \( 13 \)?}
{\( 1 \)}{\( 2 \)}{\( 3 \)}{\( 4 \)}{}{}}
\LANGes{
\Question{
¿Cuántos números impares de dos dígitos dan de resto \( 2 \) cuando se dividen por  \( 9 \) y también son divisibles por \( 13 \)?}
{\( 1 \)}{\( 2 \)}{\( 3 \)}{\( 4 \)}{}{}}
\End

\Data\ID{1003099403}
\Updated{2020-08-28 02:08:54}
\Level{B}
\SubArea{Elementary arithmetics}
\LANGen{
\Question{
The number \( x \)  when divided by \( 7 \) gives a remainder of \( 3 \). The number \( x \)  can be written in the form:}
{\( 7n+3\text{, }n\in\mathbb{N} \)}{\( 3n+7\text{, }n\in\mathbb{N} \)}{\( 7(n+3)\text{, }n\in\mathbb{N} \)}{\( 3(n+7)\text{, }n\in\mathbb{N} \)}{}{}}
\LANGpl{
\Question{
Liczba \( x \)  przy dzieleniu przez \( 7 \) daje resztę \( 3 \). Liczbę \( x \)  można więc zapisać w postaci:}
{\( 7n+3\text{, }n\in\mathbb{N} \)}{\( 3n+7\text{, }n\in\mathbb{N} \)}{\( 7(n+3)\text{, }n\in\mathbb{N} \)}{\( 3(n+7)\text{, }n\in\mathbb{N} \)}{}{}}
\LANGcs{
\Question{
Když číslo \( x \)  vydělíme \( 7 \), dostaneme zbytek \( 3 \). Číslo \( x \)  může být zapsáno ve tvaru:}
{\( 7n+3\text{, }n\in\mathbb{N} \)}{\( 3n+7\text{, }n\in\mathbb{N} \)}{\( 7(n+3)\text{, }n\in\mathbb{N} \)}{\( 3(n+7)\text{, }n\in\mathbb{N} \)}{}{}}
\LANGsk{
\Question{
Ak číslo \( x \)  vydelíme \( 7 \) dostaneme zvyšok \( 3 \). Číslo \( x \)  môžeme zapísať v tvare:}
{\( 7n+3\text{, }n\in\mathbb{N} \)}{\( 3n+7\text{, }n\in\mathbb{N} \)}{\( 7(n+3)\text{, }n\in\mathbb{N} \)}{\( 3(n+7)\text{, }n\in\mathbb{N} \)}{}{}}
\LANGes{
\Question{
Cuando el número  \( x \)  se divide por  \( 7 \) esto da de resto  \( 3 \). El número \( x \)  se puede escribir en la forma:}
{\( 7n+3\text{, }n\in\mathbb{N} \)}{\( 3n+7\text{, }n\in\mathbb{N} \)}{\( 7(n+3)\text{, }n\in\mathbb{N} \)}{\( 3(n+7)\text{, }n\in\mathbb{N} \)}{}{}}
\End

\Data\ID{1003099404}
\Updated{2020-07-15 03:07:21}
\Level{B}
\SubArea{Elementary arithmetics}
\LANGen{
\Question{
The number  \( 725233+x \) after dividing by \( 9 \) is to give a remainder of \( 5 \).  Which of the given numbers shall we substitute for \( x \)?}
{\( 1 \)}{\( 3 \)}{\( 2 \)}{\( 8 \)}{}{}}
\LANGpl{
\Question{
Liczba  \( 725233+x \) przy dzieleniu przez \( 9 \) daje resztę \( 5 \).  Liczba \( x \) może być równa?}
{\( 1 \)}{\( 3 \)}{\( 2 \)}{\( 8 \)}{}{}}
\LANGcs{
\Question{
Číslo  \( 725233+x \) po vydělení \( 9 \) má zbytek \( 5 \).  Kterým z daných čísel nahradíme \( x \)?}
{\( 1 \)}{\( 3 \)}{\( 2 \)}{\( 8 \)}{}{}}
\LANGsk{
\Question{
Číslo \( 725233+x \) po vydelení \( 9 \) dáva zvyšok \( 5 \). Ktorým z daných čísel môžme nahradiť číslo \( x \)?}
{\( 1 \)}{\( 3 \)}{\( 2 \)}{\( 8 \)}{}{}}
\LANGes{
\Question{
El número\( 725233+x \) después de ser dividido por  \( 9 \) tiene de resto  \( 5 \).  ¿Cuál de los números dados sustituiremos por \( x \)?}
{\( 1 \)}{\( 3 \)}{\( 2 \)}{\( 8 \)}{}{}}
\End

\Data\ID{1003099405}
\Updated{2020-07-15 03:07:21}
\Level{B}
\SubArea{Elementary arithmetics}
\LANGen{
\Question{
The number \( 5\cdot11\cdot17 \) has exactly:}
{eight positive integer divisors}{six positive integer divisors}{seven positive integer divisors}{five  positive integer  divisors}{}{}}
\LANGpl{
\Question{
Liczba \( 5\cdot11\cdot17 \) ma:}
{osiem dzielników naturalnych}{sześć dzielników naturalnych}{siedem dzielników naturalnych}{pięć dzielników naturalnych}{}{}}
\LANGcs{
\Question{
Číslo \( 5\cdot11\cdot17 \) má přesně:}
{osm kladných celočíselných dělitelů}{šest  kladných celočíselných dělitelů}{sedm kladných celočíselných dělitelů}{pět  kladných celočíselných dělitelů}{}{}}
\LANGsk{
\Question{
Číslo \( 5\cdot11\cdot17 \) má presne:}
{osem kladných celočíselných deliteľov}{šesť kladných celočíselných deliteľov}{sedem kladných celočíselných deliteľov}{päť kladných celočíselných deliteľov}{}{}}
\LANGes{
\Question{
El número  \( 5\cdot11\cdot17 \) tiene exactamente:}
{ocho divisores enteros positivos}{cinco divisores enteros positivos}{siete divisores enteros positivos}{cinco divisores enteros positivos}{}{}}
\End

\Data\ID{1003099406}
\Updated{2020-07-15 03:07:21}
\Level{B}
\SubArea{Elementary arithmetics}
\LANGen{
\Question{
In mathematics, the product of all positive integers less than or equal to a non-
negative integer \( n \) is denoted by \( n! \).  For example: \( 5!=5\cdot4\cdot3\cdot2\cdot1=120 \). Which of the following 
statements is true?}
{\( 16! \) is divisible by \( 91 \).}{\( 16! \) is divisible by \( 71 \).}{\( 16! \) is divisible by \( 51 \).}{\( 16! \) is divisible by \( 41 \).}{}{}}
\LANGpl{
\Question{
W matematyce iloczyn wszystkich dodatnich liczb naturalnych mniejszych lub równych liczbie naturalnej \( n \) jest oznaczony przez \( n! \). Na przykład: \( 5!=5\cdot4\cdot3\cdot2\cdot1=120 \). Które stwierdzenie jest prawdziwe?}
{\( 16! \) jest podzielna \( 91 \).}{\( 16! \) jest podzielna \( 71 \).}{\( 16! \) jest podzielna \( 51 \).}{\( 16! \) jest podzielna \( 41 \).}{}{}}
\LANGcs{
\Question{
V matematice, součin všech kladných celých čísel menších než nebo rovnající se nezápornému celému číslu \( n \) je označen \( n! \). Například: \( 5!=5\cdot4\cdot3\cdot2\cdot1=120 \). Které z následujících tvrzení je pravdivé?}
{\( 16! \) je dělitelné \( 91 \).}{\( 16! \) je dělitelné \( 71 \).}{\( 16! \) je dělitelné \( 51 \).}{\( 16! \) je dělitelné \( 41 \).}{}{}}
\LANGsk{
\Question{
V matematike sa pojmom faktoriál kladného celého čísla \( n \) označuje súčin všetkých kladných celých čísel menších alebo rovných \( n \) a označuje sa \( n! \). 
Napríklad: \( 5!=5\cdot4\cdot3\cdot2\cdot1=120 \).
Ktoré z nasledujúcich tvrdení je pravdivé?}
{\( 16! \) je deliteľné \( 91 \).}{\( 16! \) je deliteľné \( 71 \).}{\( 16! \) je deliteľné \( 51 \).}{\( 16! \) je deliteľné \( 41 \).}{}{}}
\LANGes{
\Question{
En matemáticas, el producto de todos los enteros positivos menores o iguales a un entero no negativo \( n \) se denota por \( n! \).  Por ejemplo: \( 5!=5\cdot4\cdot3\cdot2\cdot1=120 \). ¿Cuál de las siguientes afirmaciones es verdadera?}
{\( 16! \) es divisible por  \( 91 \).}{\( 16! \) es divisible por  \( 71 \).}{\( 16! \) es divisible por  \( 51 \).}{\( 16! \) es divisible por \( 41 \).}{}{}}
\End

\Data\ID{1003099407}
\Updated{2020-07-15 03:07:21}
\Level{B}
\SubArea{Elementary arithmetics}
\LANGen{
\Question{
If we convert the fraction \( \frac27 \) to a decimal, then the \( 32 \)nd digit after the decimal point is:}
{\( 8 \)}{\( 1 \)}{\( 2 \)}{\( 7 \)}{}{}}
\LANGpl{
\Question{
W rozwinięciu dziesiętnym ułamka \( \frac27 \)  na  \( 32 \)  miejscu po przecinku stoi cyfra:}
{\( 8 \)}{\( 1 \)}{\( 2 \)}{\( 7 \)}{}{}}
\LANGcs{
\Question{
Když převedeme zlomek \( \frac27 \) na desetinné číslo, pak \( 32 \). číslice za desetinnou čárkou je:}
{\( 8 \)}{\( 1 \)}{\( 2 \)}{\( 7 \)}{}{}}
\LANGsk{
\Question{
Ak premeníme zlomok \( \frac27 \) na desatinné číslo, tak \( 32. \) číslica za desatinnou čiarkou je:}
{\( 8 \)}{\( 1 \)}{\( 2 \)}{\( 7 \)}{}{}}
\LANGes{
\Question{
Si convertimos la fracción  \( \frac27 \) a decimal, entonces el dígito  \( 32 \) después del punto decimal es:}
{\( 8 \)}{\( 1 \)}{\( 2 \)}{\( 7 \)}{}{}}
\End

\Data\ID{1003118102}
\Updated{2020-07-15 03:07:21}
\Level{B}
\SubArea{Elementary arithmetics}
\LANGen{
\Question{
The number \( 10^{2010}+2 \) is divisible by:}
{\( 6 \)}{\( 5 \)}{\( 10 \)}{\( 4 \)}{}{}}
\LANGpl{
\Question{
Liczba  \( 10^{2010}+2 \) jest podzielna przez:}
{\( 6 \)}{\( 5 \)}{\( 10 \)}{\( 4 \)}{}{}}
\LANGcs{
\Question{
Číslo \( 10^{2010}+2 \) je dělitelné:}
{\( 6 \)}{\( 5 \)}{\( 10 \)}{\( 4 \)}{}{}}
\LANGsk{
\Question{
Číslo \( 10^{2010}+2 \) je deliteľné:}
{\( 6 \)}{\( 5 \)}{\( 10 \)}{\( 4 \)}{}{}}
\LANGes{
\Question{
El número \( 10^{2010}+2 \) es divisible por:}
{\( 6 \)}{\( 5 \)}{\( 10 \)}{\( 4 \)}{}{}}
\End

\Data\ID{1003118106}
\Updated{2020-07-15 03:07:21}
\Level{B}
\SubArea{Elementary arithmetics}
\LANGen{
\Question{
The number \( 11^{12}+11^{13} \) is divisible by:}
{\( 6 \)}{\( 16 \)}{\( 14 \)}{\( 5 \)}{}{}}
\LANGpl{
\Question{
Liczba \( 11^{12}+11^{13} \) jest podzielna przez:}
{\( 6 \)}{\( 16 \)}{\( 14 \)}{\( 5 \)}{}{}}
\LANGcs{
\Question{
Číslo \( 11^{12}+11^{13} \) je dělitelné:}
{\( 6 \)}{\( 16 \)}{\( 14 \)}{\( 5 \)}{}{}}
\LANGsk{
\Question{
Číslo \( 11^{12}+11^{13} \) je deliteľné:}
{\( 6 \)}{\( 16 \)}{\( 14 \)}{\( 5 \)}{}{}}
\LANGes{
\Question{
El número \( 11^{12}+11^{13} \) es divisible por:}
{\( 6 \)}{\( 16 \)}{\( 14 \)}{\( 5 \)}{}{}}
\End

\Data\ID{9000076001}
\Updated{2020-07-15 03:07:37}
\Level{B}
\SubArea{Elementary arithmetics}
\LANGen{
\Question{
In the following list identify a set of a numbers which give remainder
\(2\) after division by
\(3\), i.e. the numbers can
be written in the form \(3k + 2\),
\(k\in \mathbb{N}_{0}\).}
{\(5,\ 8,\ 11\)}{\(5,\ 10,\ 15\)}{\(3,\ 6,\ 9\)}{\(15,\ 25,\ 30\)}{\(4,\ 5,\ 6\)}{}}
\LANGpl{
\Question{
Wybierz zbiór liczb postaci
\(3k + 2\), gdzie
\(k\in \mathbb{N}_{0}\) 
(liczba, która przy dzieleniu przez
\(3\) daje
resztę \(2\)).}
{\(5,\ 8,\ 11\)}{\(5,\ 10,\ 15\)}{\(3,\ 6,\ 9\)}{\(15,\ 25,\ 30\)}{\(4,\ 5,\ 6\)}{}}
\LANGcs{
\Question{
Vyberte takovou skupinu čísel, jejíž  členy lze zapsat ve tvaru
\(3k + 2\), kde
\(k\in \mathbb{N}_{0}\) 
(jinými slovy: hledáme taková čísla, která  při  dělení číslem 
\(3\) dávají
zbytek  \(2\)).}
{\(5,\ 8,\ 11\)}{\(5,\ 10,\ 15\)}{\(3,\ 6,\ 9\)}{\(15,\ 25,\ 30\)}{\(4,\ 5,\ 6\)}{}}
\LANGsk{
\Question{
Z nasledujúcich možnosti vyberte takú množinu, v ktorej každý člen po delení \(3\) dáva zvyšok \(2\). Všetky členy množiny môžeme zapísať v tvare \(3k + 2\), kde
\(k\in \mathbb{N}_{0}\).}
{\(5,\ 8,\ 11\)}{\(5,\ 10,\ 15\)}{\(3,\ 6,\ 9\)}{\(15,\ 25,\ 30\)}{\(4,\ 5,\ 6\)}{}}
\LANGes{
\Question{
En la siguiente lista identifica un conjunto de números que dan de resto 
\(2\) después de dividirlos por \(3\), es decir, los números pueden ser escritos en la forma \(3k + 2\),
\(k\in \mathbb{N}_{0}\).}
{\(5,\ 8,\ 11\)}{\(5,\ 10,\ 15\)}{\(3,\ 6,\ 9\)}{\(15,\ 25,\ 30\)}{\(4,\ 5,\ 6\)}{}}
\End

\Data\ID{9000076002}
\Updated{2020-07-15 03:07:37}
\Level{B}
\SubArea{Elementary arithmetics}
\LANGen{
\Question{
In the following list identify a set of a numbers which give remainder
\(2\) after division by
\(5\), i.e. the numbers can
be written in the form \(5k + 2\),
\(k\in \mathbb{N}_{0}\).}
{\(37,\ 42,\ 102\)}{\(5,\ 10,\ 15\)}{\(17,\ 27,\ 100\)}{\(29,\ 47,\ 60\)}{\(41,\ 55,\ 62\)}{}}
\LANGpl{
\Question{
Wybierz zbiór liczb postaci 
 \(5k + 2\), gdzie
\(k\in \mathbb{N}_{0}\) (liczba, która przy dzieleniu przez \(5\) daje resztę \(2\)).}
{\(37,\ 42,\ 102\)}{\(5,\ 10,\ 15\)}{\(17,\ 27,\ 100\)}{\(29,\ 47,\ 60\)}{\(41,\ 55,\ 62\)}{}}
\LANGcs{
\Question{
Vyberte takovou skupinu čísel, jejíž členy lze zapsat ve tvaru
\(5k + 2\), kde
\(k\in \mathbb{N}_{0}\) 
(jinými slovy: hledáme taková čísla, která při dělení číslem
\(5\) dávají
zbytek \(2\)).}
{\(37,\ 42,\ 102\)}{\(5,\ 10,\ 15\)}{\(17,\ 27,\ 100\)}{\(29,\ 47,\ 60\)}{\(41,\ 55,\ 62\)}{}}
\LANGsk{
\Question{
Z nasledujúcich možnosti vyberte takú množinu, v ktorej každý člen po delení \(5\) dáva zvyšok \(2\). Všetky členy množiny môžeme zapísať v tvare \(5k + 2\), kde
\(k\in \mathbb{N}_{0}\).}
{\(37,\ 42,\ 102\)}{\(5,\ 10,\ 15\)}{\(17,\ 27,\ 100\)}{\(29,\ 47,\ 60\)}{\(41,\ 55,\ 62\)}{}}
\LANGes{
\Question{
En la siguiente lista identifica un conjunto de números que dan de resto
\(2\) después de dividirlos por
\(5\), es decir, los números pueden ser escritos en la forma \(5k + 2\),
\(k\in \mathbb{N}_{0}\).}
{\(37,\ 42,\ 102\)}{\(5,\ 10,\ 15\)}{\(17,\ 27,\ 100\)}{\(29,\ 47,\ 60\)}{\(41,\ 55,\ 62\)}{}}
\End

\Data\ID{9000076003}
\Updated{2020-07-15 03:07:37}
\Level{B}
\SubArea{Elementary arithmetics}
\LANGen{
\Question{
In the following list identify a set of a numbers which give remainder
\(1\) after division by
\(11\), i.e. the numbers can
be written in the form \(11k + 1\),
\(k\in \mathbb{N}_{0}\).}
{\(56,\ 122,\ 221\)}{\(21,\ 32,\ 48\)}{\(18,\ 88,\ 115\)}{\(34,\ 55,\ 70\)}{\(45,\ 56,\ 65\)}{}}
\LANGpl{
\Question{
Wskaż zbiór liczb postaci  \(11k + 1\), gdzie
\(k\in \mathbb{N}_{0}\) (liczba, która przy dzieleniu przez  \(11\)    daje resztę \(1\)).}
{\(56,\ 122,\ 221\)}{\(21,\ 32,\ 48\)}{\(18,\ 88,\ 115\)}{\(34,\ 55,\ 70\)}{\(45,\ 56,\ 65\)}{}}
\LANGcs{
\Question{
Vyberte takovou skupinu čísel, jejíž členy lze zapsat ve tvaru
\(11k + 1\), kde
\(k\in \mathbb{N}_{0}\) 
(jinými slovy: hledáme taková čísla, která při dělení číslem
\(11\) dávají
zbytek \(1\)).}
{\(56,\ 122,\ 221\)}{\(21,\ 32,\ 48\)}{\(18,\ 88,\ 115\)}{\(34,\ 55,\ 70\)}{\(45,\ 56,\ 65\)}{}}
\LANGsk{
\Question{
Z nasledujúcich možnosti vyberte takú množinu, v ktorej každý člen po delení \(11\) dáva zvyšok \(1\). Všetky členy množiny môžeme zapísať v tvare \(11k + 1\), kde
\(k\in \mathbb{N}_{0}\).}
{\(56,\ 122,\ 221\)}{\(21,\ 32,\ 48\)}{\(18,\ 88,\ 115\)}{\(34,\ 55,\ 70\)}{\(45,\ 56,\ 65\)}{}}
\LANGes{
\Question{
En la siguiente lista identifica un conjunto de números que dan de resto 
\(1\) después al dividirlos por  \(11\), es decir, los números pueden ser escritos en la forma \(11k + 1\),
\(k\in \mathbb{N}_{0}\).}
{\(56,\ 122,\ 221\)}{\(21,\ 32,\ 48\)}{\(18,\ 88,\ 115\)}{\(34,\ 55,\ 70\)}{\(45,\ 56,\ 65\)}{}}
\End

\Data\ID{9000076004}
\Updated{2020-07-15 03:07:37}
\Level{B}
\SubArea{Elementary arithmetics}
\LANGen{
\Question{
In the following list identify a set such that each element of this set is a divisor of
\(256\).}
{\(1,\ 128,\ 256\)}{\(1,\ 64,\ 123\)}{\(4,\ 8,\ 104\)}{\(1,\ 12,\ 128\)}{\(16,\ 30,\ 64\)}{}}
\LANGpl{
\Question{
Wskaż zbiór dzielników liczby
\(256\).}
{\(1,\ 128,\ 256\)}{\(1,\ 64,\ 123\)}{\(4,\ 8,\ 104\)}{\(1,\ 12,\ 128\)}{\(16,\ 30,\ 64\)}{}}
\LANGcs{
\Question{
Vyberte takovou skupinu čísel, jejíž každý člen je dělitelem čísla
\(256\).}
{\(1,\ 128,\ 256\)}{\(1,\ 64,\ 123\)}{\(4,\ 8,\ 104\)}{\(1,\ 12,\ 128\)}{\(16,\ 30,\ 64\)}{}}
\LANGsk{
\Question{
Z nasledujúcich možnosti vyberte takú množinu, v ktorej každý člen je deliteľom čísla \(256\).}
{\(1,\ 128,\ 256\)}{\(1,\ 64,\ 123\)}{\(4,\ 8,\ 104\)}{\(1,\ 12,\ 128\)}{\(16,\ 30,\ 64\)}{}}
\LANGes{
\Question{
En la siguiente lista identifica un conjunto tal que cada elemento de este conjunto sea divisor de 
\(256\).}
{\(1,\ 128,\ 256\)}{\(1,\ 64,\ 123\)}{\(4,\ 8,\ 104\)}{\(1,\ 12,\ 128\)}{\(16,\ 30,\ 64\)}{}}
\End

\Data\ID{9000076005}
\Updated{2020-07-15 03:07:37}
\Level{B}
\SubArea{Elementary arithmetics}
\LANGen{
\Question{
In the following list identify a set such that each element of this set is a divisor of
\(1\: 260\).}
{\(1,\ 36,\ 42\)}{\(4,\ 8,\ 630\)}{\(12,\ 18,\ 26\)}{\(16,\ 315,\ 1\: 260\)}{\(1,\ 17,\ 256\)}{}}
\LANGpl{
\Question{
Wskaż zbiór dzielników liczby
\(1\: 260\).}
{\(1,\ 36,\ 42\)}{\(4,\ 8,\ 630\)}{\(12,\ 18,\ 26\)}{\(16,\ 315,\ 1\: 260\)}{\(1,\ 17,\ 256\)}{}}
\LANGcs{
\Question{
Vyberte takovou skupinu čísel, jejíž každý člen je dělitelem čísla
\(1\: 260\).}
{\(1,\ 36,\ 42\)}{\(4,\ 8,\ 630\)}{\(12,\ 18,\ 26\)}{\(16,\ 315,\ 1\: 260\)}{\(1,\ 17,\ 256\)}{}}
\LANGsk{
\Question{
Z nasledujúcich možnosti vyberte takú množinu, v ktorej každý člen je deliteľom čísla
\(1\: 260\).}
{\(1,\ 36,\ 42\)}{\(4,\ 8,\ 630\)}{\(12,\ 18,\ 26\)}{\(16,\ 315,\ 1\: 260\)}{\(1,\ 17,\ 256\)}{}}
\LANGes{
\Question{
En la siguiente lista identifica un conjunto tal que cada elemento de este conjunto sea divisor de 
\(1\: 260\).}
{\(1,\ 36,\ 42\)}{\(4,\ 8,\ 630\)}{\(12,\ 18,\ 26\)}{\(16,\ 315,\ 1\: 260\)}{\(1,\ 17,\ 256\)}{}}
\End

\Data\ID{9000076006}
\Updated{2020-07-15 03:07:37}
\Level{B}
\SubArea{Elementary arithmetics}
\LANGen{
\Question{
In the following list identify a set such that each element of this set is a divisor of
\(578\).}
{\(17,\ 34,\ 289\)}{\(1,\ 2,\ 4\)}{\(13,\ 15,\ 17\)}{\(1,\ 13,\ 289\)}{\(2,\ 35,\ 578\)}{}}
\LANGpl{
\Question{
Wskaż zbiór dzielników liczby
\(578\).}
{\(17,\ 34,\ 289\)}{\(1,\ 2,\ 4\)}{\(13,\ 15,\ 17\)}{\(1,\ 13,\ 289\)}{\(2,\ 35,\ 578\)}{}}
\LANGcs{
\Question{
Vyberte takovou skupinu čísel, jejíž každý člen je dělitelem čísla
\(578\).}
{\(17,\ 34,\ 289\)}{\(1,\ 2,\ 4\)}{\(13,\ 15,\ 17\)}{\(1,\ 13,\ 289\)}{\(2,\ 35,\ 578\)}{}}
\LANGsk{
\Question{
Z nasledujúcich možnosti vyberte takú množinu, v ktorej každý člen je deliteľom čísla \(578\).}
{\(17,\ 34,\ 289\)}{\(1,\ 2,\ 4\)}{\(13,\ 15,\ 17\)}{\(1,\ 13,\ 289\)}{\(2,\ 35,\ 578\)}{}}
\LANGes{
\Question{
En la siguiente lista identifica un conjunto tal que cada elemento de este conjunto sea divisor de 
\(578\).}
{\(17,\ 34,\ 289\)}{\(1,\ 2,\ 4\)}{\(13,\ 15,\ 17\)}{\(1,\ 13,\ 289\)}{\(2,\ 35,\ 578\)}{}}
\End

\Data\ID{9000076007}
\Updated{2020-07-15 03:07:37}
\Level{B}
\SubArea{Elementary arithmetics}
\LANGen{
\Question{
Complete the statement. „The sum of any three consecutive integers ...”}
{is divisible by \(3\).}{is not divisible by \(6\).}{is divisible by \(6\).}{is not divisible by \(3\).}{is divisible by \(9\).}{}}
\LANGpl{
\Question{
Uzupełnij zdanie „Suma dowolnych trzech kolejnych liczb całkowitych ...”}
{jest podzielna przez 
\(3\).}{jest podzielna przez \(6\).}{nie jest podzielna \(6\).}{nie jest podzielna przez  \(3\).}{jest podzielna przez \(9\).}{}}
\LANGcs{
\Question{
Dokončete větu tak, aby byla pravdivá. „Součet každých tří po sobě
jdoucích celých čísel ...”}
{je dělitelný \(3\).}{není dělitelný \(6\).}{je dělitelný \(6\).}{není dělitelný \(3\).}{je dělitelný \(9\).}{}}
\LANGsk{
\Question{
Dokončite vetu tak, aby bola pravdivá. „Súčet každých troch po sebe idúcich celých čísel ...”}
{je deliteľný \(3\).}{nie je deliteľný \(6\).}{je deliteľný \(6\).}{nie je deliteľný \(3\).}{je deliteľný \(9\).}{}}
\LANGes{
\Question{
Completa la declaración. "La suma de tres enteros consecutivos ..."}
{es divisible por \(3\).}{no es divisible por \(6\).}{es divisible por \(6\).}{no es divisible por \(3\).}{es divisible por \(9\).}{}}
\End

\Data\ID{9000076008}
\Updated{2020-07-15 03:07:37}
\Level{B}
\SubArea{Elementary arithmetics}
\LANGen{
\Question{
Complete the statement. „The sum of any five consecutive integers ...”}
{is divisible by \(5\).}{is divisible by \(3\).}{is divisible by \(4\).}{is divisible by \(6\).}{is divisible by \(10\).}{}}
\LANGpl{
\Question{
Uzupełnij zdanie „Suma dowolnych pięciu  kolejnych liczb całkowitych...”}
{jest podzielna przez \(5\).}{jest podzielna przez \(3\).}{jest podzielna przez \(4\).}{jest podzielna przez \(6\).}{jest podzielna przez \(10\).}{}}
\LANGcs{
\Question{
Dokončete větu tak, aby byla pravdivá. „Součet každých pěti po sobě
jdoucích celých čísel ...”}
{je dělitelný \(5\).}{je dělitelný \(3\).}{je dělitelný \(4\).}{je dělitelný \(6\).}{je dělitelný \(10\).}{}}
\LANGsk{
\Question{
Dokončite vetu tak, aby bola pravdivá. „Súčet každých päť po sebe idúcich celých čísel ...”}
{je deliteľný \(5\).}{je deliteľný \(3\).}{je deliteľný \(4\).}{je deliteľný \(6\).}{je deliteľný \(10\).}{}}
\LANGes{
\Question{
Completa la declaración. "La suma de cinco enteros consecutivos ..."}
{es divisible por \(5\).}{es divisible por \(3\).}{es divisible por \(4\).}{es divisible por \(6\).}{es divisible por \(10\).}{}}
\End

\Data\ID{9000076009}
\Updated{2020-07-15 03:07:37}
\Level{B}
\SubArea{Elementary arithmetics}
\LANGen{
\Question{
In the following list identify a set of the numbers where all the numbers are primes.}
{\(3,\ 7,\ 89\)}{\(7,\ 15,\ 17\)}{\(8,\ 11,\ 17\)}{\(2,\ 7,\ 91\)}{\(3,\ 27,\ 81\)}{}}
\LANGpl{
\Question{
Wskaż zbiór liczb, w którym wszystkie liczby, to liczby pierwsze.}
{\(3,\ 7,\ 89\)}{\(7,\ 15,\ 17\)}{\(8,\ 11,\ 17\)}{\(2,\ 7,\ 91\)}{\(3,\ 27,\ 81\)}{}}
\LANGcs{
\Question{
Vyberte takovou skupinu čísel, jejíž každý člen má právě dva
přirozené dělitele.}
{\(3,\ 7,\ 89\)}{\(7,\ 15,\ 17\)}{\(8,\ 11,\ 17\)}{\(2,\ 7,\ 91\)}{\(3,\ 27,\ 81\)}{}}
\LANGsk{
\Question{
Z nasledujúcich možnosti vyberte takú množinu, ktorá obsahuje len prvočísla.}
{\(3,\ 7,\ 89\)}{\(7,\ 15,\ 17\)}{\(8,\ 11,\ 17\)}{\(2,\ 7,\ 91\)}{\(3,\ 27,\ 81\)}{}}
\LANGes{
\Question{
En la lista siguiente identifica un conjunto de números donde todos los números son números primos.}
{\(3,\ 7,\ 89\)}{\(7,\ 15,\ 17\)}{\(8,\ 11,\ 17\)}{\(2,\ 7,\ 91\)}{\(3,\ 27,\ 81\)}{}}
\End

\Data\ID{9000076010}
\Updated{2020-07-15 03:07:37}
\Level{B}
\SubArea{Elementary arithmetics}
\LANGen{
\Question{
In the following list identify a set of the numbers where
any number has exactly three divisors (including the number
\(1\) and
itself).}
{\(4,\ 25,\ 289\)}{\(1,\ 2,\ 3\)}{\(25,\ 36,\ 49\)}{\(1,\ 17,\ 289\)}{\(25,\ 36,\ 121\)}{}}
\LANGpl{
\Question{
Wskaż zbiór liczb, w którym każda liczba ma dokładnie trzy dzielniki (w tym liczbę  
\(1\) i
samą siebie).}
{\(4,\ 25,\ 289\)}{\(1,\ 2,\ 3\)}{\(25,\ 36,\ 49\)}{\(1,\ 17,\ 289\)}{\(25,\ 36,\ 121\)}{}}
\LANGcs{
\Question{
Vyberte takovou skupinu čísel, jejíž každý člen má právě tři
přirozené dělitele.}
{\(4,\ 25,\ 289\)}{\(1,\ 2,\ 3\)}{\(25,\ 36,\ 49\)}{\(1,\ 17,\ 289\)}{\(25,\ 36,\ 121\)}{}}
\LANGsk{
\Question{
Z nasledujúcich možnosti vyberte takú množinu, v ktorej každý člen má práve troch prirodzených deliteľov.}
{\(4,\ 25,\ 289\)}{\(1,\ 2,\ 3\)}{\(25,\ 36,\ 49\)}{\(1,\ 17,\ 289\)}{\(25,\ 36,\ 121\)}{}}
\LANGes{
\Question{
En la siguiente lista identifica un conjunto de números donde cualquier número tiene exactamente tres divisores (incluido el número
\(1\) y él mismo).}
{\(4,\ 25,\ 289\)}{\(1,\ 2,\ 3\)}{\(25,\ 36,\ 49\)}{\(1,\ 17,\ 289\)}{\(25,\ 36,\ 121\)}{}}
\End

\Data\ID{9000084901}
\Updated{2020-07-15 03:07:37}
\Level{B}
\SubArea{Elementary arithmetics}
\LANGen{
\Question{
Which of the following numbers is a prime?}
{\(17\)}{\(1\)}{\(27\)}{\(91\)}{\(289\)}{}}
\LANGpl{
\Question{
Liczbą pierwszą jest:}
{\(17\)}{\(1\)}{\(27\)}{\(91\)}{\(289\)}{}}
\LANGcs{
\Question{
Z následujících čísel vyberte prvočíslo.}
{\(17\)}{\(1\)}{\(27\)}{\(91\)}{\(289\)}{}}
\LANGsk{
\Question{
Ktoré z nasledujúcich čísel je prvočíslo?}
{\(17\)}{\(1\)}{\(27\)}{\(91\)}{\(289\)}{}}
\LANGes{
\Question{
¿Cuál de los siguientes números es número primo?}
{\(17\)}{\(1\)}{\(27\)}{\(91\)}{\(289\)}{}}
\End

\Data\ID{9000084902}
\Updated{2020-07-15 03:07:37}
\Level{B}
\SubArea{Elementary arithmetics}
\LANGen{
\Question{
In the following list find the set which does not contain any prime number.}
{\(91,\ 243\)}{\(13,\ 100\)}{\(2,\ 4\)}{\(29,\ 81\)}{\(101,\ 211\)}{}}
\LANGpl{
\Question{
Wskaż zbiór liczb zawierający liczby złożone.}
{\(91,\ 243\)}{\(13,\ 100\)}{\(2,\ 4\)}{\(29,\ 81\)}{\(101,\ 211\)}{}}
\LANGcs{
\Question{
Z následujících skupin čísel vyberte tu, která neobsahuje žádné
prvočíslo.}
{\(91,\ 243\)}{\(13,\ 100\)}{\(2,\ 4\)}{\(29,\ 81\)}{\(101,\ 211\)}{}}
\LANGsk{
\Question{
Z nasledujúcich možností vyberte takú, ktorá neobsahuje žiadne prvočísla.}
{\(91,\ 243\)}{\(13,\ 100\)}{\(2,\ 4\)}{\(29,\ 81\)}{\(101,\ 211\)}{}}
\LANGes{
\Question{
En la siguiente lista encuentra el conjunto que no contiene ningún número primo.}
{\(91,\ 243\)}{\(13,\ 100\)}{\(2,\ 4\)}{\(29,\ 81\)}{\(101,\ 211\)}{}}
\End

\Data\ID{9000084903}
\Updated{2020-07-15 03:07:37}
\Level{B}
\SubArea{Elementary arithmetics}
\LANGen{
\Question{
In the following list find the set which contains only prime numbers.}
{\(13,\ 131\)}{\(1,\ 31,\ 211\)}{\(289,\ 291\)}{\(17,\ 169\)}{\(51,\ 97\)}{}}
\LANGpl{
\Question{
Wskaż zbiór, który zawiera tylko liczby pierwsze.}
{\(13,\ 131\)}{\(1,\ 31,\ 211\)}{\(289,\ 291\)}{\(17,\ 169\)}{\(51,\ 97\)}{}}
\LANGcs{
\Question{
Z následujících skupin čísel vyberte tu, která obsahuje jen prvočísla.}
{\(13,\ 131\)}{\(1,\ 31,\ 211\)}{\(289,\ 291\)}{\(17,\ 169\)}{\(51,\ 97\)}{}}
\LANGsk{
\Question{
Z nasledujúcich možností vyberte takú, ktorá obsahuje len prvočísla.}
{\(13,\ 131\)}{\(1,\ 31,\ 211\)}{\(289,\ 291\)}{\(17,\ 169\)}{\(51,\ 97\)}{}}
\LANGes{
\Question{
En la siguiente lista encuentra el conjunto que  no contiene ningún número primo.}
{\(13,\ 131\)}{\(1,\ 31,\ 211\)}{\(289,\ 291\)}{\(17,\ 169\)}{\(51,\ 97\)}{}}
\End

\Data\ID{9000084904}
\Updated{2020-07-15 03:07:37}
\Level{B}
\SubArea{Elementary arithmetics}
\LANGen{
\Question{
From the following list find the number which has just three proper divisors.}
{\(49\)}{\(21\)}{\(75\)}{\(100\)}{\(250\)}{}}
\LANGpl{
\Question{
Wskaż liczbę, która ma tylko trzy dzielniki.}
{\(49\)}{\(21\)}{\(75\)}{\(100\)}{\(250\)}{}}
\LANGcs{
\Question{
Z následujících čísel vyberte to, které má právě tři kladné dělitele.}
{\(49\)}{\(21\)}{\(75\)}{\(100\)}{\(250\)}{}}
\LANGsk{
\Question{
Z nasledujúcich čísel vyberte také, ktoré má práve tri kladné delitele.}
{\(49\)}{\(21\)}{\(75\)}{\(100\)}{\(250\)}{}}
\LANGes{
\Question{
En la siguiente lista encuentra el número que tiene solo tres divisores propios.}
{\(49\)}{\(21\)}{\(75\)}{\(100\)}{\(250\)}{}}
\End

\Data\ID{9000084905}
\Updated{2020-07-15 03:07:37}
\Level{B}
\SubArea{Elementary arithmetics}
\LANGen{
\Question{
In the following list find the number such that the prime factorization of this
number contains just two different primes or powers of two different primes.}
{\(100\)}{\(5\)}{\(25\)}{\(120\)}{\(121\)}{}}
\LANGpl{
\Question{
Wskaż liczbę, która w rozkładzie na czynniki pierwsze ma dokładnie dwie różne liczby pierwsze.}
{\(100\)}{\(5\)}{\(25\)}{\(120\)}{\(121\)}{}}
\LANGcs{
\Question{
Z následujících čísel vyberte to, které má v prvočíselném rozkladu
právě dvě různá prvočísla.}
{\(100\)}{\(5\)}{\(25\)}{\(120\)}{\(121\)}{}}
\LANGsk{
\Question{
Z následujúcich čísel vyberte také, ktoré má v prvočíselnom rozklade práve dve rôzne prvočísla.}
{\(100\)}{\(5\)}{\(25\)}{\(120\)}{\(121\)}{}}
\LANGes{
\Question{
En la siguiente lista encuentra el número cuya factorización en factores primos contenga solo dos primos diferentes o potencias de dos primos diferentes.}
{\(100\)}{\(5\)}{\(25\)}{\(120\)}{\(121\)}{}}
\End

\Data\ID{9000084906}
\Updated{2020-07-15 03:07:37}
\Level{B}
\SubArea{Elementary arithmetics}
\LANGen{
\Question{
In the following list find the number such that the prime factorization of this
number contains exactly one cube power.}
{\(24\)}{\(12\)}{\(63\)}{\(196\)}{\(420\)}{}}
\LANGpl{
\Question{
Wskaż liczbę, w której w rozkładzie na czynniki pierwsze występuje jedna liczba pierwsza w potędze trzeciej.}
{\(24\)}{\(12\)}{\(63\)}{\(196\)}{\(420\)}{}}
\LANGcs{
\Question{
Z následujících čísel vyberte to, které má v prvočíselném rozkladu
právě jedno prvočíslo ve třetí mocnině.}
{\(24\)}{\(12\)}{\(63\)}{\(196\)}{\(420\)}{}}
\LANGsk{
\Question{
Z následujúcich čísel vyberte také, ktoré má v prvočíselnom rozklade práve jedno prvočíslo v tretej mocnine.}
{\(24\)}{\(12\)}{\(63\)}{\(196\)}{\(420\)}{}}
\LANGes{
\Question{
En la siguiente lista encuentra el número cuya factorización en factores primos contenga exactamente una potencia al cubo.}
{\(24\)}{\(12\)}{\(63\)}{\(196\)}{\(420\)}{}}
\End

\Data\ID{9000084907}
\Updated{2020-09-21 12:09:45}
\Level{B}
\SubArea{Elementary arithmetics}
\LANGen{
\Question{
Among the following numbers, find the one that has the greatest number of different primes in its prime factorization.}
{\(330\)}{\(21\)}{\(100\)}{\(486\)}{\(1\: 024\)}{}}
\LANGpl{
\Question{
Wskaż liczbę, która w rozkładzie na czynniki pierwsze ma najwięcej liczb pierwszych.}
{\(330\)}{\(21\)}{\(100\)}{\(486\)}{\(1\: 024\)}{}}
\LANGcs{
\Question{
Z následujících čísel vyberte to, které má v prvočíselném rozkladu
nejvíce různých prvočísel.}
{\(330\)}{\(21\)}{\(100\)}{\(486\)}{\(1\: 024\)}{}}
\LANGsk{
\Question{
Z následujúcich čísel vyberte také, ktoré má v prvočíselnom rozklade najviac rôznych prvočísel.}
{\(330\)}{\(21\)}{\(100\)}{\(486\)}{\(1\: 024\)}{}}
\LANGes{
\Question{
En la siguiente lista encuentra el número cuya factorización en factores primos contenga más números primos diferentes en comparación con la factorización  de otros números de esta lista.}
{\(330\)}{\(21\)}{\(100\)}{\(486\)}{\(1\: 024\)}{}}
\End

\Data\ID{9000084908}
\Updated{2021-01-08 02:01:39}
\Level{B}
\SubArea{Elementary arithmetics}
\LANGen{
\Question{
From the following list of numbers choose the one  that has in its  prime factorization the highest power of a prime.}
{\(1\: 024\)}{\(21\)}{\(100\)}{\(330\)}{\(486\)}{}}
\LANGpl{
\Question{
Wskaż liczbę, która ma najwyższą potęgę  w rozkładzie na czynniki pierwsze.}
{\(1\: 024\)}{\(21\)}{\(100\)}{\(330\)}{\(486\)}{}}
\LANGcs{
\Question{
Z následujících čísel vyberte to, které v prvočíselném rozkladu
obsahuje prvočíslo v nejvyšší mocnině.}
{\(1\: 024\)}{\(21\)}{\(100\)}{\(330\)}{\(486\)}{}}
\LANGsk{
\Question{
Z následujúcich čísel vyberte také, ktoré obsahuje v prvočíselnom rozklade prvočíslo v najvyššej mocnine.}
{\(1\: 024\)}{\(21\)}{\(100\)}{\(330\)}{\(486\)}{}}
\LANGes{
\Question{
Entre los números en la siguiente lista encuentra los números que tienen mayor potencia en la factorización en factores primos  que los otros números de la lista.}
{\(1\: 024\)}{\(21\)}{\(100\)}{\(330\)}{\(486\)}{}}
\End

\Data\ID{9000084909}
\Updated{2020-09-21 12:09:10}
\Level{B}
\SubArea{Elementary arithmetics}
\LANGen{
\Question{
Among the following numbers, find the number so that in its factorization into primes each prime is squared.}
{\(36\)}{\(24\)}{\(120\)}{\(360\)}{\(512\)}{}}
\LANGpl{
\Question{
Wskaż liczbę, która w rozkładzie na czynniki pierwsze nie zawiera potęgi innej niż  2.}
{\(36\)}{\(24\)}{\(120\)}{\(360\)}{\(512\)}{}}
\LANGcs{
\Question{
Z následujících čísel vyberte to, které má v prvočíselném rozkladu
prvočísla pouze ve druhé mocnině.}
{\(36\)}{\(24\)}{\(120\)}{\(360\)}{\(512\)}{}}
\LANGsk{
\Question{
Z následujúcich čísel vyberte také, ktoré má v prvočíselnom rozklade prvočísla práve v druhej mocnine.}
{\(36\)}{\(24\)}{\(120\)}{\(360\)}{\(512\)}{}}
\LANGes{
\Question{
Entre los números en la siguiente lista encuentra el número cuya factorización en factores primos no contenga una potencia diferente de 2.}
{\(36\)}{\(24\)}{\(120\)}{\(360\)}{\(512\)}{}}
\End

\Data\ID{9000084910}
\Updated{2020-07-15 03:07:37}
\Level{B}
\SubArea{Elementary arithmetics}
\LANGen{
\Question{
In the following list find the number such that the prime factorization of this
number contains only one prime.}
{\(125\)}{\(15\)}{\(100\)}{\(250\)}{\(768\)}{}}
\LANGpl{
\Question{
Wskaż liczbę, która w rozkładzie na czynniki pierwsze ma tylko jedną liczbę pierwszą.}
{\(125\)}{\(15\)}{\(100\)}{\(250\)}{\(768\)}{}}
\LANGcs{
\Question{
Z následujících čísel vyberte to, které nemá v prvočíselném rozkladu
různá prvočísla.}
{\(125\)}{\(15\)}{\(100\)}{\(250\)}{\(768\)}{}}
\LANGsk{
\Question{
Z nasledujúcich čísel vyberte také, ktoré nemá v prvočíselnom rozklade rôzne prvočísla.}
{\(125\)}{\(15\)}{\(100\)}{\(250\)}{\(768\)}{}}
\LANGes{
\Question{
En la siguiente lista encuentra el número cuya factorización en factores primos contenga solo un número primo.}
{\(125\)}{\(15\)}{\(100\)}{\(250\)}{\(768\)}{}}
\End

\Data\ID{9000085605}
\Updated{2020-07-15 03:07:37}
\Level{B}
\SubArea{Elementary arithmetics}
\LANGen{
\Question{
Find the product of all one-digit primes and round to nearest hundreds.}
{\(200\)}{\(100\)}{\(300\)}{\(400\)}{}{}}
\LANGpl{
\Question{
Znajdź iloczyn wszystkich jednocyfrowych liczb pierwszych i zaokrągli do setki.}
{\(200\)}{\(100\)}{\(300\)}{\(400\)}{}{}}
\LANGcs{
\Question{
Součin všech jednociferných prvočísel zaokrouhlený na stovky je roven:}
{\(200\)}{\(100\)}{\(300\)}{\(400\)}{}{}}
\LANGsk{
\Question{
Súčin všetkých jednociferných prvočísel zaokrúhlený na stovky sa rovná:}
{\(200\)}{\(100\)}{\(300\)}{\(400\)}{}{}}
\LANGes{
\Question{
Encuentra el producto de todos los números primos de un dígito y redondea a las centenas más cercanas.}
{\(200\)}{\(100\)}{\(300\)}{\(400\)}{}{}}
\End

\Data\ID{9000085606}
\Updated{2020-07-15 03:07:37}
\Level{B}
\SubArea{Elementary arithmetics}
\LANGen{
\Question{
Find the product of all divisors of \(12\) 
and round to nearest hundreds.}
{\(1\: 700\)}{\(1\: 200\)}{\(600\)}{\(100\)}{}{}}
\LANGpl{
\Question{
Znajdź iloczyn wszystkich dzielników liczby \(12\) i
zaokrągli do setki.}
{\(1\: 700\)}{\(1\: 200\)}{\(600\)}{\(100\)}{}{}}
\LANGcs{
\Question{
Součin všech dělitelů čísla \(12\) 
zaokrouhlený na stovky je roven:}
{\(1\: 700\)}{\(1\: 200\)}{\(600\)}{\(100\)}{}{}}
\LANGsk{
\Question{
Súčin všetkých deliteľov čísla \(12\) 
zaokrúhlený na stovky sa rovná:}
{\(1\: 700\)}{\(1\: 200\)}{\(600\)}{\(100\)}{}{}}
\LANGes{
\Question{
Encuentra el producto de todos los divisores de  \(12\) 
y redondea a las centenas más cercanas.}
{\(1\: 700\)}{\(1\: 200\)}{\(600\)}{\(100\)}{}{}}
\End

\Data\ID{9000115601}
\Updated{2020-07-15 03:07:37}
\Level{B}
\SubArea{Elementary arithmetics}
\LANGen{
\Question{
Complete the following statement: „The number is divisible by two if and only if ...”}
{the last digit of this number is even.}{the sum of its digits is divisible by two.}{the sum of its digits is even.}{the last digit of this number is \(2,\ 3,\ 6\) 
or \(8\).}{}{}}
\LANGpl{
\Question{
Uzupełnij zdanie „Liczba jest podzielna przez dwa wtedy i tylko wtedy, gdy ...”}
{ostatnia cyfra tej liczby jest parzysta.}{suma jej cyfr jest podzielna przez dwa.}{suma jej cyfr jest parzysta.}{ostatnia cyfra tej liczby to  \(2,\ 3,\ 6\) 
lub \(8\).}{}{}}
\LANGcs{
\Question{
Číslo je dělitelné dvěma, je-li}
{poslední číslice sudá.}{ciferný součet dělitelný dvěma.}{ciferný součet sudý.}{poslední číslice \(2,\ 3,\ 6\) 
nebo \(8\).}{}{}}
\LANGsk{
\Question{
Prirodzené číslo je deliteľné číslom dva práve vtedy, ak je}
{jeho posledná číslica párna.}{číslom dva deliteľný jeho ciferný súčet.}{jeho ciferný súčet párny.}{jeho posledná číslica \(2,\ 3,\ 6\) 
alebo \(8\).}{}{}}
\LANGes{
\Question{
Completa la siguiente proposición lógica: "El número es divisible por dos si y solo si ..."}
{el último dígito de este número es par.}{la suma de sus dígitos es divisible por dos.}{la suma de sus dígitos es par.}{el último dígito de este número es \(2,\ 3,\ 6\) 
u \(8\).}{}{}}
\End

\Data\ID{9000115602}
\Updated{2020-07-15 03:07:37}
\Level{B}
\SubArea{Elementary arithmetics}
\LANGen{
\Question{
Complete the following statement: „The number is divisible by three if and only if ...”}
{the sum of its digits is divisible by three.}{the number constituted from the last two digits is divisible by three.}{the sum of its digits is odd.}{the last digit of this number is \(3,\ 6\) 
or \(9\).}{}{}}
\LANGpl{
\Question{
Uzupełnij zdanie „Liczba jest podzielna przez trzy wtedy i tylko wtedy, gdy ...”}
{suma jej cyfr jest podzielna przez trzy.}{jej dwie ostatnie cyfry tworzą liczbę podzielną przez trzy.}{suma jej cyfr jest nieparzysta.}{ostatnia cyfra tej liczby to  \(3,\ 6\) 
lub \(9\).}{}{}}
\LANGcs{
\Question{
Číslo je dělitelné třemi, je-li}
{ciferný součet dělitelný třemi.}{poslední dvojčíslí dělitelné třemi.}{ciferný součet lichý.}{poslední číslice \(3,\ 6\) 
nebo \(9\).}{}{}}
\LANGsk{
\Question{
Prirodzené číslo je deliteľné číslom tri práve vtedy, ak je}
{číslom tri deliteľný jeho ciferný súčet.}{číslom tri deliteľné jeho posledné dvojčíslie.}{jeho ciferný súčet nepárny.}{jeho posledná číslica \(3,\ 6\) 
alebo \(9\).}{}{}}
\LANGes{
\Question{
Completa la siguiente proposición lógica: "El número es divisible por tres si y solo si ..."}
{la suma de sus dígitos es divisible por tres.}{el número formado por los dos últimos dígitos es divisible por tres.}{la suma de sus dígitos es impar.}{el último dígito de este número es  \(3,\ 6\) 
o \(9\).}{}{}}
\End

\Data\ID{9000115603}
\Updated{2020-07-15 03:07:37}
\Level{B}
\SubArea{Elementary arithmetics}
\LANGen{
\Question{
Complete the following statement: „The number is divisible by four if and only if ...”}
{the number constituted from the last two digits is divisible by four.}{the sum of its digits is divisible by four.}{the last digit of this number is \(4\).}{the last digit of this number is even.}{}{}}
\LANGpl{
\Question{
Uzupełnij zdanie „Liczba jest podzielna przez cztery wtedy i tylko wtedy, gdy ...”}
{jej dwie ostatnie cyfry tworzą liczbę podzielną przez  cztery.}{suma jej cyfr jest podzielna przez cztery.}{ostatnia cyfra tej liczby to 
 \(4\).}{ostatnia cyfra tej liczby jest parzysta.}{}{}}
\LANGcs{
\Question{
Číslo je dělitelné čtyřmi, je-li}
{poslední dvojčíslí dělitelné čtyřmi.}{ciferný součet dělitelný čtyřmi.}{poslední číslice čtyřka.}{poslední číslice sudá.}{}{}}
\LANGsk{
\Question{
Prirodzené číslo je deliteľné číslom štyri práve vtedy, ak je}
{číslom štyri deliteľné jeho posledné dvojčíslie.}{číslom štyri deliteľný jeho ciferný súčet.}{jeho posledná číslica štyri.}{jeho posledná číslica párna.}{}{}}
\LANGes{
\Question{
Completa la siguiente proposición lógica: "El número es divisible por cuatro si y solo si ..."}
{el número formado por los dos últimos dígitos es divisible por cuatro.}{la suma de sus dígitos es divisible por cuatro.}{el último dígito de este número es  \(4\).}{el último dígito de este número es par.}{}{}}
\End

\Data\ID{9000115604}
\Updated{2020-07-15 03:07:37}
\Level{B}
\SubArea{Elementary arithmetics}
\LANGen{
\Question{
Complete the following statement: „The number is divisible by five if and only if ...”}
{the last digit of this number is \(5\) 
or \(0\).}{the sum of its digits is divisible by five.}{it is divisible by two and three.}{the last digit of this number is odd.}{}{}}
\LANGpl{
\Question{
Uzupełnij zdanie „Liczba jest podzielna przez pięć wtedy i tylko wtedy, gdy ...”}
{ostatnia cyfra tej liczby to 
 \(5\) 
lub \(0\).}{suma jej cyfr jest podzielna przez pięć.}{jest podzielna przez dwa i trzy.}{ostatnia cyfra tej liczby jest nieparzysta.}{}{}}
\LANGcs{
\Question{
Číslo je dělitelné pěti, je-li}
{poslední číslice pět nebo nula.}{ciferný součet dělitelný pěti.}{dělitelné dvěma a třemi současně.}{poslední číslice lichá.}{}{}}
\LANGsk{
\Question{
Prirodzené číslo je deliteľné číslom päť práve vtedy, ak je}
{jeho posledná číslica nula alebo päť.}{číslom päť deliteľný jeho ciferný súčet.}{deliteľné číslami dva a tri.}{jeho posledná číslica nepárna.}{}{}}
\LANGes{
\Question{
Completa la siguiente proposición lógica: "El número es divisible por cinco si y solo si ..."}
{el último dígito de este número es \(5\) 
o \(0\).}{la suma de sus dígitos es divisible por cinco.}{es divisible por dos y tres.}{el último dígito de este número es impar.}{}{}}
\End

\Data\ID{9000115605}
\Updated{2020-07-15 03:07:37}
\Level{B}
\SubArea{Elementary arithmetics}
\LANGen{
\Question{
Complete the following statement: „The number is divisible by six if and only if ...”}
{it is divisible by both two and three.}{the sum of its digits is divisible by both two and three.}{the sum of its digits is even and the last digit of this number is
\(3\).}{the last digit of this number is \(6\).}{}{}}
\LANGpl{
\Question{
Uzupełnij zdanie „Liczba jest podzielna przez sześć  wtedy i tylko wtedy, gdy ...”}
{jest podzielna dwa i trzy.}{suma jej cyfr jest podzielna przez dwa i trzy.}{suma jej cyfr jest parzysta, a ostatnia cyfra tej liczby to 
\(3\).}{ostatnia cyfra tej liczby to  \(6\).}{}{}}
\LANGcs{
\Question{
Číslo je dělitelné šesti, je-li}
{dělitelné dvěma a třemi současně.}{ciferný součet dělitelný dvěma a současně třemi.}{ciferný součet sudý a poslední cifra je 3.}{jeho poslední číslice šest.}{}{}}
\LANGsk{
\Question{
Prirodzené číslo je deliteľné číslom šesť práve vtedy, ak je}
{deliteľné číslami dva a tri.}{číslom dva a tri deliteľný jeho ciferný súčet.}{jeho ciferný súčet párny a jeho posledná číslica je tri.}{jeho posledná číslica šesť.}{}{}}
\LANGes{
\Question{
Completa la siguiente proposición lógica: "El número es divisible por seis si y solo si ..."}
{es divisible por dos y por tres.}{la suma de sus dígitos es divisible por dos y por tres.}{la suma de sus dígitos es par y el último dígito de este número es 
\(3\).}{el último dígito de este número es \(6\).}{}{}}
\End

\Data\ID{9000115606}
\Updated{2020-07-15 03:07:37}
\Level{B}
\SubArea{Elementary arithmetics}
\LANGen{
\Question{
Complete the following statement: „The number is divisible by eight if and only if ...”}
{the number constituted from the last three digits is divisible by eight.}{the sum of its digits is divisible by eight.}{it is divisible by both two and four.}{the number constituted from the last two digits is divisible by eight.}{}{}}
\LANGpl{
\Question{
Uzupełnij zdanie „Liczba jest podzielna przez osiem wtedy i tylko wtedy, gdy ...”}
{jej trzy ostatnie cyfry tworzą liczbę podzielną przez osiem.}{suma jej cyfr jest podzielna przez osiem.}{jest podzielna przez dwa i cztery.}{jej dwie ostatnie cyfry tworzą liczbę podzielną przez  osiem.}{}{}}
\LANGcs{
\Question{
Číslo je dělitelné osmi, je-li}
{poslední trojčíslí dělitelné osmi.}{ciferný součet dělitelný osmi.}{dělitelné dvěma a čtyřmi současně.}{poslední dvojčíslí dělitelné osmi.}{}{}}
\LANGsk{
\Question{
Prirodzené číslo je deliteľné číslom osem práve vtedy, ak je}
{číslom osem deliteľné jeho posledné trojčíslie.}{číslom osem deliteľný jeho ciferný súčet.}{deliteľné číslami dva a štyri.}{číslom osem deliteľné jeho posledné dvojčíslie.}{}{}}
\LANGes{
\Question{
Completa la siguiente proposición lógica: "El número es divisible por ocho si y solo si ..."}
{el número formado por los últimos tres dígitos es divisible por ocho.}{la suma de sus dígitos es divisible por ocho.}{es divisible por dos y cuatro a la vez.}{el número constituido por los dos últimos dígitos es divisible por ocho.}{}{}}
\End

\Data\ID{9000115607}
\Updated{2020-07-15 03:07:37}
\Level{B}
\SubArea{Elementary arithmetics}
\LANGen{
\Question{
Complete the following statement: „The number is divisible by nine if and only if ...”}
{the sum of its digits is divisible by nine.}{the number constituted from the last two digits is divisible by nine.}{the sum of its digits is odd.}{the last digit of this number is \(9\).}{}{}}
\LANGpl{
\Question{
Uzupełnij zdanie „Liczba jest podzielna przez dziewięć  wtedy i tylko wtedy, gdy ...”}
{suma jej cyfr jest podzielna przez dziewięć.}{jej dwie ostatnie cyfry tworzą liczbę podzielną przez dziewięć.}{suma jej cyfr jest nieparzysta.}{ostatnia cyfra tej liczby to  \(9\).}{}{}}
\LANGcs{
\Question{
Číslo je dělitelné devíti, je-li}
{ciferný součet dělitelný devíti.}{poslední dvojčíslí dělitelné devíti.}{ciferný součet lichý.}{poslední číslice devět.}{}{}}
\LANGsk{
\Question{
Prirodzené číslo je deliteľné číslom deväť práve vtedy, ak je}
{číslom deväť deliteľný jeho ciferný súčet.}{číslom deväť deliteľné jeho posledné dvojčíslie.}{jeho ciferný súčet nepárny.}{posledná číslica devät.}{}{}}
\LANGes{
\Question{
Completa la siguiente proposición lógica: "El número es divisible por nueve si y solo si ..."}
{la suma de sus dígitos es divisible por nueve.}{el número formado por los dos últimos dígitos es divisible por nueve.}{la suma de sus dígitos es impar.}{el último dígito de este número es \(9\).}{}{}}
\End

\Data\ID{9000115608}
\Updated{2020-07-15 03:07:37}
\Level{B}
\SubArea{Elementary arithmetics}
\LANGen{
\Question{
Complete the following statement: „The number is divisible by ten if and only if ...”}
{the last digit of this number is \(0\).}{the sum of its digits is divisible by ten.}{the number constituted from the last two digits is divisible by five.}{the last digit of this number is even.}{}{}}
\LANGpl{
\Question{
Uzupełnij zdanie „Liczba jest podzielna przez dziesięć  wtedy i tylko wtedy, gdy ...”}
{ostatnia cyfra tej liczby to 
 \(0\).}{suma jej cyfr jest podzielna przez dziesięć.}{jej dwie ostatnie cyfry tworzą liczbę podzielną przez pięć.}{ostatnia jej cyfra jest parzysta.}{}{}}
\LANGcs{
\Question{
Číslo je dělitelné deseti, je-li}
{poslední číslice nula.}{ciferný součet dělitelný deseti.}{poslední dvojčíslí dělitelné pěti.}{poslední číslice sudá.}{}{}}
\LANGsk{
\Question{
Prirodzené číslo je deliteľné číslom desať práve vtedy, ak je}
{jeho posledná číslica nula.}{číslom desať deliteľný jeho ciferný súčet.}{číslom päť deliteľné jeho posledné dvojčíslie.}{jeho posledná číslica párna.}{}{}}
\LANGes{
\Question{
Completa la siguiente proposición lógica: "El número es divisible por diez si y solo si ..."}
{el último dígito de este número es \(0\).}{la suma de sus dígitos es divisible por diez.}{el número formado por los dos últimos dígitos es divisible por cinco.}{el último dígito de este número es par.}{}{}}
\End

\Data\ID{9000115609}
\Updated{2020-09-21 12:09:55}
\Level{B}
\SubArea{Elementary arithmetics}
\LANGen{
\Question{
Complete the following statement: „A number is divisible by twelve if and only if ...”}
{it is divisible by three and four.}{the sum of its digits is divisible by both two and three.}{the sum of the digits is even and the last digit of this number is odd.}{the sum of the digits is odd and the last digit of this number is even.}{}{}}
\LANGpl{
\Question{
Uzupełnij zdanie „Liczba jest podzielna przez dwanaście   wtedy i tylko wtedy, gdy ...”}
{jest podzielna przez trzy i cztery.}{suma jej cyfr jest podzielna przez dwa i trzy.}{suma jej cyfr jest parzysta, a ostatnia cyfra tej liczby jest nieparzysta.}{suma jej cyfr jest nieparzysta, a ostatnia cyfra tej liczby jest parzysta.}{}{}}
\LANGcs{
\Question{
Číslo je dělitelné dvanácti, je-li}
{dělitelné současně třemi a čtyřmi.}{ciferný součet dělitelný dvěma a třemi současně.}{ciferný součet sudý a poslední dvojčíslí je liché.}{poslední číslice sudá a ciferný součet je lichý.}{}{}}
\LANGsk{
\Question{
Prirodzené číslo je deliteľné číslom dvanásť práve vtedy, ak je}
{deliteľné číslami tri a štyri.}{číslom tri a štyri deliteľný jeho ciferný súčet.}{jeho ciferný súčet párny a posledné dvojčíslie nepárne.}{jeho ciferný súčet nepárny a posledné dvojčíslie párne.}{}{}}
\LANGes{
\Question{
Completa la siguiente proposición lógica: "El número es divisible por doce si y solo si ..."}
{divisible por tres y por cuatro.}{la suma de sus dígitos es divisible por dos y por tres.}{la suma de los dígitos es par y el último dígito de este número es impar.}{la suma de los dígitos es impar y el último dígito de este número es par.}{}{}}
\End

\Data\ID{9000115610}
\Updated{2020-09-21 12:09:01}
\Level{B}
\SubArea{Elementary arithmetics}
\LANGen{
\Question{
Complete the following statement: „A number is divisible by fifteen if and only if ...”}
{it is divisible by both three and five.}{the sum of its digits is divisible by both three and five.}{the sum of its digits is odd and divisible by five.}{the last digit of this number is either \(5\) 
or \(0\).}{}{}}
\LANGpl{
\Question{
Uzupełnij zdanie „Liczba jest podzielna przez piętnaście  wtedy i tylko wtedy, gdy ...”}
{jest podzielna przez trzy i pięć.}{suma jej cyfr jest podzielna przez trzy i pięć.}{suma jej cyfr jest nieparzysta i podzielna przez pięć.}{ostatnia cyfra tej liczby to  \(5\) 
lub \(0\).}{}{}}
\LANGcs{
\Question{
Číslo je dělitelné patnácti, je-li}
{dělitelné současně třemi a pěti.}{ciferný součet dělitelný třemi a pěti současně.}{ciferný součet lichý a dělitelný pěti.}{poslední číslice pět nebo nula.}{}{}}
\LANGsk{
\Question{
Prirodzené číslo je deliteľné číslom pätnásť práve vtedy, ak je}
{deliteľné číslami tri a päť.}{číslom tri a päť deliteľný jeho ciferný súčet.}{jeho ciferný súčet párny a deliteľný piatimi.}{jeho posledná číslica päť alebo nula.}{}{}}
\LANGes{
\Question{
Completa la siguiente proposición lógica: "El número es divisible por quince si y solo si ..."}
{divisible por tres y por cinco.}{la suma de sus dígitos es divisible por tres y cinco a la vez.}{la suma de sus dígitos es impar y divisible por cinco.}{el último dígito de este número es  \(5\) 
o \(0\).}{}{}}
\End
